%% 第31回編集委員会（物理第5巻）：組 原W8
%% p5-15.tex  原-15　原子炉の中の中性子：核分裂と原子炉（★2〜5・3段：N-1）
%% 中身は計画書 tools/phys5-plan.md の B「原-15」。公式の持ち主は A（A8・A10・A37・A78〜A83）。
%%   ほかの章は計画書 B の節の番号で指した（原-11 ⑤⑥・例題10、原-12 ③⑧、原-13 ⑥⑧・例題10、原-16）。
%% 数値はすべて Python で計算した：checks/物理第5巻/原-15/calc_p5_15.py（出力 calc_p5_15_out.txt）、
%%   code_p5_15.py（⑦⑧の囲みのプログラム。出力 code_p5_15_out.txt）、kanji_p5_15.py（常用漢字の確かめ）。
%% 前の巻・数学の本との約束（計画書 A-6・C-6-2）
%%   ・Q6：前付の保存則の地図「中性子：生まれる数－吸収される数、中性子のバランス、増倍率」→ ④（地図ではもれを「出る量」に入れた、と書いた）
%%   ・Q18：熱-16:136・熱-5 ⑦（原子力の蒸気は低温なので効率が低い）→ ③⑬
%%   ・Q20：熱-16:335（ITER 50 MW → 500 MW、Q＝10）→ ⑫・例題13 のじぶんでとこう
%%   ・Q21・Q25：熱-16:341 と道しるべ → ③（原子炉のしくみ）、使い終わった燃料は原-16
%%   ・Q24：ph2-back:385「ウランの濃縮」→ 熱-14 例題12 の 1.0043 → ⑤・例題5（約340段）
%%   ・Q60：電-16 の原子力発電 → ③⑬
%%   ・M10：数学 9-9:1358（臨界は極がちょうど0、制御棒で極の場所を調整）→ ④⑧
%%   ・M11・M12：数学 9-13・9-back の表（原-15 ← 9-9）→ ⑧
%%   ・M17：8-12 ⑦（球のラプラシアン、S6）→ ⑦。M20：6-2 のもっと知りたい（ルンゲ・クッタ、S13）→ ⑧・例題9
%%   ・6-14（つながった2つの量、⑦消去法）→ ⑧⑩。6-4（指数で増える・減る、⑤時定数）→ ④⑥⑩⑪
%%   ・計画書 C-4 の2：200 MeV・0.0065（原-11・15・16）、崩壊熱の割合（原-15・16）をそろえる。原-16 の見本の計算（3年運転）も同じ式。
%% 記号（この章。計画書 C-1）：増倍率 k（k_∞・k_eff。この章で波数・ばね定数・クーロンの定数は使わない）、
%%   反応度 ρ（この章で密度を記号にしない。原子の数密度は N）、遅発中性子の割合 β（原-11 の v/c、β線とは別）、
%%   1回の分裂で出る中性子の数 ν（N-5。光の振動数は f）、中性子束 φ・断面積 σ・巨視的断面積 Σ（原-13 と同じ）、
%%   拡散係数 D、移動面積 M^2、バックリング B^2（原-11 の結合エネルギー B とは章が別）、中性子の世代の時間 ℓ（1点炉の式では Λ）、
%%   先行核の崩壊定数 λ（原-12 と同じ。第4部）、先行核の数 C、中性子の数（出力に比例）n、4因子 η・ε・p・f、
%%   出力 P、崩壊熱の割合 P/P_0（原-16 と同じ）、運転した時間 T_0、エネルギー E、1回の反応のエネルギー Q（原-11 の Q 値）。
%% 出典（11月11日に作業環境から確かめた。Wikipedia は使っていない）
%%   ・核分裂の発見：ノーベル化学賞 1944年 ハーン「重い原子核の分裂の発見」（https://www.nobelprize.org/prizes/chemistry/1944/hahn/facts/）。
%%       同じ所の伝記：「1938年の終わり」にシュトラスマンと分裂を見つけ、最初の論文は 1939年1月6日と2月10日の Naturwissenschaften。
%%       マイトナーとフリッシュが理論で説明したこと（同じ facts のページ）。
%%   ・フェルミ：ノーベル物理学賞 1938年（遅い中性子による核反応）、最初の制御された連鎖反応は 1942年12月2日、シカゴのスタジアムの下の
%%       スカッシュのコート（https://www.nobelprize.org/prizes/physics/1938/fermi/biographical/）。
%%   ・1回の分裂のエネルギーの内訳（ウラン235：破片の運動 166.2、即発中性子 4.8、即発γ 8.0、破片のβ 7.0、破片のγ 7.2、ニュートリノ 9.6 MeV）：
%%       A. Algora ら, "Beta-decay studies for applied and basic nuclear physics", Eur. Phys. J. A (2021), arXiv:2007.07918 の表1（もとは Madland 2006）。
%%       同じ論文：運転中の炉の熱の約7 %が分裂生成物のβ崩壊から、止めたあとはこの崩壊熱が主になり、取り除けないと
%%       福島第一のような事故になる。崩壊熱を最初に見積もったのは Way と Wigner。
%%   ・熱中性子（0.0253 eV）の断面積：B. Pritychenko, arXiv:2503.18990（2025）の表2・表3（ENDF/B-VIII.1・JEFF-3.3・JENDL-5.0 ほか）。
%%       ウラン235 の分裂 586.1・584.5・586.7 b（JENDL-5.0 は 586.7）、捕獲 99.4・99.6・99.4 b、ウラン238 の捕獲 2.683 b、
%%       水素1 の捕獲 0.3326 b、キセノン135 の捕獲 2.664×10⁶・2.654×10⁶・2.777×10⁶ b。本文は 585 b・99 b・2.68 b・約 2.6×10⁶ b に丸めた。
%%       一次資料で確かめ直した（出典の組、11月11日）：IAEA Neutron Data Standards 2017 の表11（熱中性子定数、0.0253 eV）
%%       https://nds.iaea.org/standards/std2017/Standards2017_TNC.txt を curl で読んだ：ウラン235 の分裂 587.3±1.4 b・捕獲 99.5±1.3 b。
%%       評価ずみのライブラリ（上の 584.5〜586.7 b）と 0.5 % 以内で、本文の 585 b・99 b はそのままにした（η は 2.08 のまま）。
%%       NIST NCNR の散乱長と断面積の表（Sears 1992。https://www.ncnr.nist.gov/resources/n-lengths/list.html を curl で読んだ）：
%%       ¹H の吸収 0.3326 b、²³⁸U 2.68 b、²³⁵U の吸収（分裂＋捕獲）680.9 b。キセノン135 はこの表にない（上の Pritychenko の表で確かめた。
%%       同じ表の Atlas〔Mughabghab〕の欄は 2.65±0.11×10⁶ b）。
%%   ・ν＝2.43（ウラン235・熱中性子）：IAEA Neutron Data Standards 2017 の表11（上と同じファイル）の
%%       "Nu-bar, Total average neutron yield per 1 fission" 2.425±0.011 を丸めた（出典の組が 11月11日に curl で読んだ）。JEFF-3.1 の評価は 2.4367。
%%   ・遅発中性子の割合 β＝0.0065（ウラン235）：米国エネルギー省 DOE-HDBK-1019/1-93（Nuclear Physics and Reactor Theory）の6群の表の合計 0.0065
%%       （WebSearch の要約）と、Akyurek の論文 arXiv:2101.09338 の式 β＝ν_d/ν（JEFF-3.1.1 と ENDF/B-VIII.0 から ν_d＝1.57×10⁻²）で照らし合わせた。
%%       6群の半減期と割合（54.51 s・0.038 … 0.179 s・0.026）は arXiv:2101.09338 の表2.1。これから1群の λ＝0.0785 /s（平均 12.7 s）を計算した
%%       （DOE の手引きは平均 12.5 s）。本文は λ＝0.08 /s。臭素87 の 55.7 s も同じ論文。
%%       〔要確認：ν_d の一次資料〕（出典の組、11月11日）：ν＝2.425（IAEA の標準）で割ると 1.57×10⁻²÷2.425＝0.00647 で 0.0065 と合う。
%%       ν_d そのものを評価ずみのファイルで読もうとしたが、IAEA の nds.iaea.org/relnsd/delayedn/ は先行核ごとの P_n の表で ν_d の合計がなく、
%%       INDC の報告（nds.iaea.org/publications/indc/）は 403、OECD/NEA の JANIS web は表を出すのにブラウザのセッションが要り読めなかった。
%%       www.osti.gov・www.standards.doe.gov（DOE の手引き）は応答がなかった。
%%   ・中性子の世代の時間：DOE-HDBK-1019 は即発中性子の世代の時間を約 5×10⁻⁵ s とする（WebSearch の要約）。本文は「10⁻⁴ s ほど」とし、
%%       例題では ℓ＝1×10⁻⁴ s「とします」（計画書と同じ）。
%%   ・ヨウ素135・キセノン135：半減期 6.58±0.03 h・9.14±0.02 h（ENSDF。IAEA LiveChart の API
%%       https://nds.iaea.org/relnsd/v1/data?fields=ground_states&nuclides=135i を curl で読んだ。出典の組、11月11日）。
%%       はじめは二次資料の 6.57 h を使っていたので、6.58 h に直し、例題11・図・じぶんでとこうを計算し直した（calc_p5_15.py）。
%%       熱中性子によるウラン235 の分裂の収率（IAEA LiveChart の API、fields=cumulative_fy・independent_fy&parents=235u）：
%%       ヨウ素135 の累積収率 6.39 %（テルル135 からの分もふくむ）、キセノン135 が分裂で直接できる分（独立収率。基底状態 0.069 %
%%       ＋準安定状態 0.167 %）約0.24 %、キセノン135 の累積収率 6.61 %（約96 % がヨウ素135 から）。はじめは DOE-HDBK-1019 の要約の
%%       0.3 % を使っていたので、本文を「約0.2 %」に、例題11 を γ_I＝0.064・γ_X＝0.0024 に直した。中性子束 3×10¹³ /(cm²·s) は例題の仮の値。
%%   ・ウェイ・ウィグナーの式（崩壊熱の目安）：係数は本によって 0.062〜0.066 と違う。本文は計画書と同じ 0.066 で、目安の式と書いた。
%%       停止直後に定格の 6〜7 %（JAEA の解説の要約：WebSearch）と、Algora らの約7 %で照らし合わせた。〔要確認：係数の一次資料〕
%%       （出典の組、11月11日：今回届くようになった一次資料の中に、この式の係数をのせたものは見つからなかった。）
%%   ・軽水炉の燃料（ウラン235 3〜5 %）、福島第一の事故の経過（1〜3号機は自動停止、外部電源の喪失、津波で非常用ディーゼル発電機などが被水し
%%       6号機を除き交流電源を失った、注水の停止 1号機 約14時間・2・3号機 約6時間半、燃料の溶融、ジルコニウムと水蒸気の反応で水素、
%%       1・3・4号機の水素爆発、INES レベル7）：環境省「放射線による健康影響等に関する統一的な基礎資料（令和7年度版）」HTML 版
%%       02-02-01・02-02-03・06-01-02・06-01-03・06-01-04・06-02-01（https://www.env.go.jp/chemi/rhm/current/）。
%%   ・核融合：ローソンの条件（D-T の点火に n τ_E ≥ 1.6×10²⁰ m⁻³·s、T≈26 keV。理想の点火温度 4.3 keV。JET の Q_sci 0.64）：
%%       S. E. Wurzel, S. C. Hsu, arXiv:2105.10954 の表II と本文。NIF の 2022年12月5日の実験（レーザー 2.05 MJ、核融合 3.15 MJ、利得 1.5）：
%%       Phys. Rev. Lett. 132, 065102 (2024) の要旨と、Kritcher ら Phys. Rev. E 109, 025204 (2024) の要旨（WebSearch。二次資料2つで照らし合わせた）。
%%       ITER（50 MW → 500 MW、Q＝10）は熱-16 と同じ数（熱-16 の出典）。
%%   ・原子の質量（D・T・⁴He・n・⁹²Kr・¹⁴¹Ba・²³⁵U）と存在比 0.7204 %：計画書 C-3（NIST、AME2020）。炭素の燃焼熱 393.5 kJ/mol：熱-16・ph2-back と同じ。
%%       水の蒸発熱 100 ℃で 2257 kJ/kg：熱-1 と同じ（IAPWS-IF97）。
%% 〔仮〕（例題のために置いた値。本文に「とします」）：ε＝1.03・p＝0.75・f＝0.90、k_∞＝1.10・M²＝60 cm²、温度係数 −3 pcm/K、
%%   中性子束 3×10¹³ /(cm²·s)、海水の温度上昇 7 K、⑦のプログラムの D＝3.0 cm・Σa＝0.05 /cm。
\makeatletter
% 中性子（小さな玉）  \txpn{(x,y)}
\newcommand\txpn[1]{\filldraw[fill=bkBlue!25, draw=bkBlue!80!black, line width=0.6pt] #1 circle (1.3);}
% 原子核（大きな玉）  \txpnuc{(x,y)}{半径}{色}
\newcommand\txpnuc[3]{\shade[ball color=#3!55] #1 circle (#2); \draw[bkInk, line width=0.5pt] #1 circle (#2);}
\tikzset{txp lab/.style={font=\sffamily\footnotesize, text=bkInk}, txp arr/.style={line width=1pt, -{Stealth[length=2.4mm]}},
  txp dim/.style={bkGray, line width=0.5pt, {Stealth[length=1.8mm]}-{Stealth[length=1.8mm]}},
  txp n/.style={bkBlue!80!black, line width=0.9pt, -{Stealth[length=2mm]}}}
\makeatother

\phchapter{15}{原子炉の中の中性子：\\核分裂と原子炉}{原子炉は、どうやって反応を一定に保つ？}{\bklevelrange{2}{5}}{原子炉の中の中性子：核分裂と原子炉}
\bkmichishirube{6-2,6-4,6-14,8-12,8-13,9-9,熱-8,熱-14,原-11,原-12,原-13}{原-16}

\begin{goal}
\begin{itemize}[label={\color{bkTheme}●}, leftmargin=1.6\zw, itemsep=0pt, topsep=1pt]
\item \bklevel{2}\enspace ウラン235 の核分裂で、1回に約 200 MeV のエネルギーと2〜3個の中性子が出ることと、連鎖反応・原子炉のしくみを、図で説明できる。
\item \bklevel{3}\enspace 増倍率 $k$・反応度・遅発中性子を使って、原子炉の出力の増え方（炉の周期）を見積もれる。
\item \bklevel{5}\enspace 4因子公式と1群の拡散方程式で、増倍率と、原子炉が臨界になる大きさを求められる。
\item \bklevel{5}\enspace 1点炉の動特性の式・反応度係数・キセノンの式を、微分方程式として読める。
\item \bklevel{3}\enspace 崩壊熱と核融合のエネルギーを数で見積もり、原子炉を止めたあとも冷やし続けなければならないわけを説明できる。
\end{itemize}
\end{goal}

\phkoushiki{%
\phk{原子力などから電気のエネルギーを得る}{中学理科}{ph:原-15:sec3}
\phk{原子力は核分裂でエネルギーを出す}{高校物理基礎}{ph:原-15:sec1}
\phk{核分裂と連鎖反応（臨界 $k=1$）}{高校物理}{ph:原-15:sec2}
}

\begin{hatena}
\noindent\begin{minipage}[c]{0.55\linewidth}\raggedright
原子炉の中では、ウランの原子核が分かれるたびに中性子が飛び出し、その中性子が、次の原子核を分けます。1回分かれるごとに中性子が2個以上出るのなら、反応はねずみ算のように増えていきそうです。

でも、原子力発電所は、何か月も同じ出力で電気をつくり続けます。原子炉は、どうやって、この反応を増えも減りもしないように保っているのでしょう。
\end{minipage}\hfill
\begin{minipage}[c]{0.42\linewidth}\centering
\begin{tikzpicture}[x=1mm, y=1mm]
  \fill[bkBlue!10] (2,2) rectangle (50,38);
  \draw[bkInk, line width=1.1pt] (2,2) rectangle (50,38);
  \foreach \x in {8,14,20,26,32,38,44} {\fill[bkOrange!60] (\x-1.2,6) rectangle (\x+1.2,30); \draw[bkInk, line width=0.4pt] (\x-1.2,6) rectangle (\x+1.2,30);}
  \foreach \x in {11,23,35,47} {\fill[bkGray!70] (\x-1,20) rectangle (\x+1,44); \draw[bkInk, line width=0.4pt] (\x-1,20) rectangle (\x+1,44);}
  \txpn{(17,13)} \txpn{(29,24)} \txpn{(41,11)}
  \draw[txp n] (17,13) -- (21.5,16.5); \draw[txp n] (29,24) -- (31,19); \draw[txp n] (41,11) -- (36.5,14);
  \node[txp lab, anchor=south] at (27,45.5) {制御棒（上下に出し入れする）};
  \node[txp lab, anchor=north] at (26,1) {燃料と水（中性子は丸）};
\end{tikzpicture}
\end{minipage}
\end{hatena}

\noindent 原子の中心にある原子核には、とても大きなエネルギーがしまわれています。原-11 では、重い原子核が2つに分かれると、質量の一部がエネルギーに変わることを見ました。この章では、そのエネルギーを、とぎれずに、しかも一定の量だけ取り出す装置である\textbf{原子炉}を調べます。\bkidx{げんしろ}{原子炉}主役は、目に見えない小さな粒、\textbf{中性子}です。中性子が「何個生まれ、何個吸われ、何個外へにげるか」を数えることが、原子炉の物理のすべての出発点になります。

\begin{yattemiyou}[ドミノで連鎖をまねしよう]
\textbf{用意するもの}：ドミノ60個くらい（消しゴムなどでもよい）、ストップウォッチ\par\vspace{2pt}
\noindent\begin{minipage}[c]{0.64\linewidth}
\begin{enumerate}[label={\bkcircled[bkGreen]{\arabic*}}, leftmargin=2.4\zw, itemsep=1pt, topsep=0pt]
\item 人が通らない机の上などに、ドミノを1列に15個並べる（1個が1個を倒す：$k=1$ のまね）。
\item 別の所に、1個が2個を倒すように枝分かれさせて並べる（1段め1個、2段め2個、3段め4個、4段め8個：$k=2$ のまね）。
\item どちらも、はしの1個を倒し、倒れる数の増え方と、全部が倒れるまでの時間をはかる。
\item 1列の途中に1か所だけ間をあけて、①をもう一度倒す（$k<1$ のまね）。
\end{enumerate}
\end{minipage}\hfill
\begin{minipage}[c]{0.33\linewidth}\centering
\begin{tikzpicture}[x=1mm, y=1mm]
  \foreach \i in {0,...,5} {\fill[bkOrange!70] (2+\i*4.4,40) rectangle (3.2+\i*4.4,45);}
  \node[txp lab, anchor=west] at (28,42.5) {$k=1$};
  \fill[bkOrange!70] (2,15) rectangle (3.2,20);
  \foreach \y in {7,25} {\fill[bkOrange!70] (8,\y) rectangle (9.2,\y+5);}
  \foreach \y in {2.5,11.5,20.5,29.5} {\fill[bkOrange!70] (14,\y) rectangle (15.2,\y+4.5);}
  \foreach \y in {0,4.5,9,13.5,18,22.5,27,31.5} {\fill[bkOrange!70] (20,\y) rectangle (21.2,\y+3.5);}
  \node[txp lab, anchor=west] at (24,17.5) {$k=2$};
\end{tikzpicture}
\end{minipage}
\end{yattemiyou}

\noindent やってみようの答え：1列（$k=1$）では、ドミノはいつも1個ずつ、同じ速さで倒れ続けます。枝分かれ（$k=2$）では、倒れる数が1、2、4、8と倍々に増え、4段めまでの15個が、1列の15個よりずっと短い時間で倒れ終わります。途中に間をあけた列（$k<1$）では、そこで止まります。原子炉の中の中性子も、1個が次の世代に平均何個の中性子を残すか（増倍率 $k$）で、増えるか・一定か・減るかが決まります。ただし、ドミノは1回倒れたら終わりですが、原子炉では、新しい燃料の原子核がいつも次の中性子を待っています。だから、$k$ がほんの少しでも1より大きいと、増え方は止まりません。

\begin{omoidasou}
\bkrecall{原-11}{原子核の1核子あたりの結合エネルギーは、鉄やニッケルのあたりがいちばん大きい（約 8.8 MeV）。ウランのような重い原子核（約 7.6 MeV）は、2つに分かれると固く結びついた原子核になり、質量の差 $\Delta m$ が $\Delta mc^2$ のエネルギーになる（原-11 例題3：$^{235}$U＋n → $^{141}$Ba＋$^{92}$Kr＋3n で \qty{173}{MeV}）。}
\bkrecall{原-12}{放射性の原子核は、1個1個が偶然にこわれ、数は $\dfrac{\dd N}{\dd t}=-\lambda N$ で減る。半減期は $T_{1/2}=\ln2/\lambda$。親がこわれて子ができ、子もこわれる「親と子」の式（原-12 ⑧）。}
\bkrecall{原-13}{中性子は電気を持たないので、原子核にぶつかるまでまっすぐ進む。1個の原子核に当たる「的の大きさ」が断面積 $\sigma$（単位 b＝$\qty{e-28}{m^2}$）、数密度 $N$ をかけた $\Sigma=N\sigma$ が巨視的断面積（1 m 進むあいだに当たる回数）。中性子の流れの強さは中性子束 $\phi$。水素と何回か衝突すると、中性子は約 \qty{0.025}{eV} の熱中性子になる（2 MeV から約18回：原-13 例題10）。}
\bkrecall{熱-14}{ウランの濃縮では、軽い $^{235}$UF$_6$ が重い $^{238}$UF$_6$ の約1.0043倍の速さで小さな穴からもれ出す（グレアムの法則、熱-14 例題12）。}
\bkrecall{熱-8}{発電所では、熱で水を水蒸気にしてタービンを回し、水蒸気を海水で冷やして水にもどす（ランキンサイクル）。}
\bkrecall{6-4}{$\dfrac{\dd y}{\dd t}=ay$ の答えは $y=y_0e^{at}$。$a>0$ なら増え、$a<0$ なら減る。$1/|a|$ が時定数。}
\bkrecall{6-14}{2つの量がたがいに影響しながら変わるときは、2つの式を同じ時刻の値で同時に進める。1つの式にまとめると、答えは2つの指数の和になる（6-14⑦）。}
\end{omoidasou}

\noindent この章の記号：増倍率は $k$ と書きます（この章では、波数やばね定数には使いません）。反応度は $\rho$（ロー）と書くので、この章では密度を記号にしません。$\beta$（ベータ）は遅発中性子の割合で、原-11 の $v/c$ や β線とは別のものです。$\nu$（ニュー）は\textbf{1回の分裂で出る中性子の数}です（高校の教科書は光の振動数を $\nu$ と書きますが、この本では光の振動数を $f$ と書きます：原-2）。

%% ======================== 物理の道具 ========================
\bksection{doguu}{物理の道具}

\phstage{me}{核分裂と連鎖反応を、図で見る}{2}

\bksubsection{1}{核分裂：ウランの原子核が2つに分かれる}
\label{ph:原-15:sec1}
1938年の終わり、ドイツのハーンとシュトラスマンは、ウランに中性子を当てると、ウランよりずっと軽い\textbf{バリウム}ができることを、化学の分析で見つけました。最初の論文は1939年のはじめに出ています。同じ研究を長く続けてきたマイトナーと、そのおいのフリッシュは、これを「ウランの原子核が、ほぼ半分の大きさの2つの原子核に分かれた」と考えて説明しました。原子核が2つに分かれることを\textbf{核分裂}といいます。\bkidx{かくぶんれつ}{核分裂}ハーンは、この発見で1944年のノーベル化学賞を受けました。

\begin{bkfigure}
\begin{tikzpicture}[x=0.84mm, y=0.9mm]
  % 中性子が来る
  \txpn{(4,20)} \draw[txp n] (5.5,20) -- (14,20);
  \node[txp lab, anchor=north] at (6,17) {中性子};
  % ウラン235
  \txpnuc{(24,20)}{8}{bkGreen}
  \node[txp lab] at (24,20) {$^{235}$U};
  \draw[bkGray, txp arr] (34,20) -- (42,20);
  % ゆれる原子核
  \begin{scope}[shift={(56,20)}]
    \fill[bkGreen!45] (-5,-3) rectangle (5,3);
    \shade[ball color=bkGreen!55] (-6,0) circle (6.2); \shade[ball color=bkGreen!55] (6,0) circle (5.6);
    \fill[bkGreen!45] (-1.5,-2.6) rectangle (1.5,2.6);
  \end{scope}
  \node[txp lab, anchor=north, align=center] at (56,10) {くびれて\\ゆれる};
  \draw[bkGray, txp arr] (70,20) -- (78,20);
  % 2つの破片と中性子
  \txpnuc{(90,30)}{5.4}{bkOrange}  \node[txp lab] at (90,30) {Ba};
  \txpnuc{(90,9)}{4.7}{bkPurple}   \node[txp lab] at (90,9) {Kr};
  \draw[bkOrange!80!black, txp arr] (95,33) -- (103,38);
  \draw[bkPurple!80!black, txp arr] (95,6) -- (103,1);
  \txpn{(104,20)} \txpn{(111,25)} \txpn{(111,15)}
  \draw[txp n] (105.5,20) -- (113,20); \draw[txp n] (112.4,25.7) -- (119,28.5); \draw[txp n] (112.4,14.3) -- (119,11.5);
  \node[txp lab, anchor=west, align=left] at (114,20) {中性子\\2〜3個};
  \node[txp lab, anchor=south, text=bkRed!80!black] at (90,37) {約 200 MeV};
\end{tikzpicture}
\bkfigcaption{ウラン235 の核分裂の一例。中性子を吸ったウランの原子核は、くびれて大きくゆれ、2つの原子核（分裂片）に分かれる。分裂片は大きな速さで飛び出し、同時に中性子が2〜3個出る。分かれ方は1通りではなく、いろいろな組になる（図はバリウムとクリプトンの組）。}\label{ph:原-15:fig-fission}
\end{bkfigure}

なぜエネルギーが出るのでしょう。原-11 ⑤で、1核子（陽子と中性子）あたりの結合エネルギーの曲線を見ました。図\ref{ph:原-15:fig-BA}のように、ウラン235 は1核子あたり約 \qty{7.6}{MeV} で結びついています。分かれてできるバリウム141・クリプトン92 は約 8.3〜\qty{8.5}{MeV} で、1核子あたり約 \qty{0.8}{MeV} 固く結びつきます。核子は全部で約230個あるので、
\[ 0.8\ \text{MeV}\times230\approx180\ \text{MeV} \]
ほどの結合エネルギーが余り、分裂片の運動エネルギーなどになって出てきます。原-11 例題3 で、質量の差から正しく計算すると \qty{173}{MeV} でした。

\begin{bkfigure}
\begin{tikzpicture}
\begin{axis}[bkplot, width=106mm, height=56mm, xmin=0, xmax=250, ymin=6.5, ymax=9.3, xtick={50,100,150,200,250}, ytick={7,8,9}, xlabel={質量数}, ylabel={1核子あたりの結合エネルギー（MeV）}]
  \draw[bkGray, line width=1pt, dashed] plot[smooth] coordinates {(axis cs:12,7.68) (axis cs:16,7.98) (axis cs:56,8.79) (axis cs:92,8.51) (axis cs:141,8.33) (axis cs:235,7.59)};
  \fill[bkBlue!80!black] (axis cs:4,7.07) circle (1.8pt); \node[font=\sffamily\footnotesize, anchor=west, inner sep=2pt] at (axis cs:5,7.0) {$^{4}$He};
  \fill[bkBlue!80!black] (axis cs:12,7.68) circle (1.8pt); \node[font=\sffamily\footnotesize, anchor=north west, inner sep=2pt] at (axis cs:12,7.68) {$^{12}$C};
  \fill[bkBlue!80!black] (axis cs:16,7.98) circle (1.8pt); \node[font=\sffamily\footnotesize, anchor=east, inner sep=0.5pt] at (axis cs:15.5,8.17) {$^{16}$O};
  \fill[bkBlue!80!black] (axis cs:56,8.79) circle (1.8pt); \node[font=\sffamily\footnotesize, anchor=south, inner sep=2pt] at (axis cs:56,8.79) {$^{56}$Fe};
  \fill[bkBlue!80!black] (axis cs:92,8.51) circle (1.8pt); \node[font=\sffamily\footnotesize, anchor=south, inner sep=2pt] at (axis cs:92,8.51) {$^{92}$Kr};
  \fill[bkBlue!80!black] (axis cs:141,8.33) circle (1.8pt); \node[font=\sffamily\footnotesize, anchor=south, inner sep=2pt] at (axis cs:141,8.33) {$^{141}$Ba};
  \fill[bkBlue!80!black] (axis cs:235,7.59) circle (1.8pt); \node[font=\sffamily\footnotesize, anchor=south, inner sep=2pt] at (axis cs:235,7.59) {$^{235}$U};
  \draw[bkRed!80!black, line width=1pt, -{Stealth[length=2.4mm]}] (axis cs:228,7.45) to[bend left=20] (axis cs:145,8.2);
  \draw[bkRed!80!black, line width=1pt, -{Stealth[length=2.4mm]}] (axis cs:228,7.45) to[bend left=25] (axis cs:96,8.38);
  \node[font=\sffamily\footnotesize, text=bkRed!80!black, anchor=south, align=center] at (axis cs:165,6.6) {分かれると、1核子あたり\\約 0.8 MeV 固くなる};
\end{axis}
\end{tikzpicture}
\bkfigcaption{1核子あたりの結合エネルギー（点は NIST の原子の質量から計算した値。破線は点をなめらかにつないだもの。たての軸は 6.5 MeV から。重水素 $^2$H は 1.11 MeV で、ずっと下）。ウランのような重い原子核は、2つに分かれると、曲線のいちばん高い所（鉄）に近づく。}\label{ph:原-15:fig-BA}
\end{bkfigure}

\phboxhead{1回の分裂で出るエネルギーの行き先}
分裂した直後のエネルギーだけでなく、そのあとの分裂片の放射性崩壊（β崩壊とγ線）でもエネルギーが出ます。ウラン235 の1回の分裂のエネルギーの内訳は、表\ref{ph:原-15:tab-energy}のとおりです。ニュートリノ（原-11 ⑪）は、ほとんど何にも当たらずに原子炉の外へ出ていくので、熱にはなりません。そのほかに、中性子が燃料や水に吸われるときにもγ線が出て熱になります。そこで、原子炉の熱になるエネルギーは、\textbf{1回の分裂あたり約 \qty{200}{MeV}} と見積もります。\bkidx{かくぶんれつのえねるぎー}{核分裂のエネルギー}

\begin{bkcenter}
\begin{bktable}{colspec={Q[l,m,wd=52mm] Q[r,m,wd=20mm] Q[l,m,wd=50mm]}, rowsep=1pt, row{1}={bg=bkTheme!10}, cells={font=\small}}
エネルギーの行き先 & MeV & いつ出るか \\
分裂片の運動エネルギー & 166.2 & 分裂した瞬間 \\
即発中性子の運動エネルギー & 4.8 & 分裂した瞬間 \\
即発γ線 & 8.0 & 分裂した瞬間 \\
分裂片のβ崩壊（電子） & 7.0 & あとから（秒〜年） \\
分裂片のγ線 & 7.2 & あとから（秒〜年） \\
ニュートリノ（炉の外へ） & 9.6 & あとから（秒〜年） \\
合計 & 202.8 & \\
\end{bktable}
\bktabcaption{ウラン235 の1回の熱中性子による分裂で出るエネルギーの内訳（Algora ら、2021年の論文の表）。「あとから」の β と γ の合計 \qty{14.2}{MeV} は、熱になる分（約 \qty{193}{MeV}）の約7\,\%。これが、⑪の崩壊熱のもとになる。}\label{ph:原-15:tab-energy}
\end{bkcenter}

\begin{phmijika}[化学の反応と、けたがいくつ違う？]
石炭（炭素）が燃えるとき、炭素原子1個あたりに出るエネルギーは約 \qty{4.1}{eV} です（燃焼熱 \qty{393.5}{kJ/mol} を、1 mol の原子の数でわった値）。核分裂の \qty{200}{MeV}＝$\qty{2e8}{eV}$ は、その約5千万倍です。化学の反応は原子のまわりの電子の結びつきが変わるだけですが、核分裂は、原子核の中の、ずっと強い力の結びつきが変わるからです。だから、原子力発電所は、同じ電気をつくるのに、火力発電所の数百万分の1の質量の燃料ですみます（例題1）。そのかわり、燃え残りではなく、放射線を出す分裂片が燃料の中にたまっていきます（⑪、原-16）。
\end{phmijika}

\phboxhead{遅い中性子が、分裂を起こしやすい}
1回の分裂では、中性子が2個のときも3個のときもあり、平均すると\textbf{2.43個}出ます。この数を $\nu$（ニュー）と書きます。\bkidx{いっかいのぶんれつででるちゅうせいしのかず}{1回の分裂で出る中性子の数}出たばかりの中性子は、平均約 \qty{2}{MeV} の速い中性子です。ウラン235 は、遅い中性子（熱中性子、約 \qty{0.025}{eV}）が当たると、とてもよく分裂します。分裂の断面積は約 \qty{585}{b} で、原子核の見かけの大きさ（数 b）の100倍以上です。一方、天然のウランの99.3\,\%をしめるウラン238 は、熱中性子では分裂せず、ときどき中性子を吸う（捕獲断面積 \qty{2.68}{b}）だけです。だから、ふつうの原子炉では、出たばかりの速い中性子を、水などにぶつけて遅くしてから使います（③）。

\bksubsection{2}{連鎖反応と増倍率}
\label{ph:原-15:sec2}
1回の分裂で出た中性子のうち、何個かが別のウラン235 に吸われて、次の分裂を起こします。その分裂で出た中性子が、また次の分裂を起こします。このように、反応が次の反応を起こしてつながっていくことを\textbf{連鎖反応}といいます。\bkidx{れんさはんのう}{連鎖反応}

ある「世代」の中性子が、次の世代に平均何個の中性子を残すかを、\textbf{増倍率} $k$ といいます。\bkidx{ぞうばいりつ}{増倍率}
\[ k=\frac{\text{次の世代の中性子の数}}{\text{今の世代の中性子の数}} \]
1回の分裂で2.43個出ても、全部が次の分裂を起こすわけではありません。燃料以外の物（水・構造の金属・制御棒）に吸われたり、原子炉の外へにげたりするからです。$k$ は、それを全部差し引いた「生き残って次の分裂を起こす割合」で決まります。

\begin{bkfigure}
\begin{tikzpicture}[x=1mm, y=1mm]
  % k=1
  \node[txp lab, anchor=west] at (0,46) {$k=1$（臨界）};
  \foreach \i in {0,...,4} {\txpnuc{(6+\i*13,36)}{3.4}{bkGreen} \ifnum\i<4 \draw[txp n] (9.6+\i*13,36) -- (15.4+\i*13,36);\fi}
  \foreach \i in {0,...,4} {\draw[bkGray, line width=0.6pt, -{Stealth[length=1.6mm]}] (6+\i*13,32.6) -- (8+\i*13,28);}
  \node[txp lab, anchor=north west, align=left] at (0,26.5) {残りの中性子は、吸われるか、\\外へにげる};
  % k=2
  \begin{scope}[shift={(68,0)}]
  \node[txp lab, anchor=west] at (0,46) {$k=2$（超臨界）};
  \txpnuc{(5,24)}{3.2}{bkGreen}
  \foreach \y in {14,34} {\draw[txp n] (8.2,24) -- (14,\y); \txpnuc{(17,\y)}{3.2}{bkGreen}}
  \foreach \y in {9,19,29,39} {\txpnuc{(30,\y)}{2.8}{bkGreen}}
  \foreach \a/\b in {14/9,14/19,34/29,34/39} {\draw[txp n] (20,\a) -- (27,\b);}
  \foreach \y in {6,11,16,21,26,31,36,41} {\txpnuc{(43,\y)}{2.0}{bkGreen}}
  \foreach \a/\b in {9/6,9/11,19/16,19/21,29/26,29/31,39/36,39/41} {\draw[bkBlue!80!black, line width=0.7pt, -{Stealth[length=1.5mm]}] (32.8,\a) -- (40.8,\b);}
  \node[txp lab, anchor=north] at (24,4) {1→2→4→8};
  \end{scope}
\end{tikzpicture}
\bkfigcaption{増倍率 $k$ と連鎖反応。左：1回の分裂から、ちょうど1個の中性子が次の分裂を起こすと、分裂の数は一定（臨界）。右：2個ずつ起こすと、世代ごとに2倍になる（超臨界）。}\label{ph:原-15:fig-chain}
\end{bkfigure}

\begin{itemize}[leftmargin=1.6\zw, itemsep=1pt, topsep=2pt]
\item $k>1$：中性子が世代ごとに増える。\textbf{超臨界}。
\item $k=1$：中性子の数が一定。\textbf{臨界}。\bkidx{りんかい}{臨界}一定の出力で運転している原子炉は、いつもこの状態です。
\item $k<1$：中性子が世代ごとに減り、連鎖反応はやがて止まる。\textbf{未臨界}。
\end{itemize}
「臨界」は、ニュースでは「危ない状態」のように聞こえることがありますが、物理のことばとしては「中性子の数がつり合って、一定の出力で落ち着いている状態」のことです。

はじめの世代の中性子の数を $n_0$ とすると、$g$ 世代あとの数は、$k$ を $g$ 回かけて
\pheq{n_g=n_0\,k^g}{個 ＝ 個 × 1}
です（3-11 の倍々でふえる形）。

\begin{phyomi}[$n_g=n_0k^g$]
「中性子の数は、1世代ごとに $k$ 倍になる」と読みます。$k$ が1よりほんの少し大きいだけでも、世代を重ねると大きく増えます。$k=1.001$ なら、100世代で $1.001^{100}=1.105$ 倍、1万世代で約2万倍です。原子炉の中の1世代は、10$^{-4}$ 秒（1万分の1秒）ほどしかありません（⑥）。だから、$k$ を1にぴったり合わせ続けることが、原子炉の運転のいちばん大事な仕事になります。
\end{phyomi}

\begin{bkfigure}
\begin{tikzpicture}
\begin{axis}[bkplot, width=112mm, height=66mm, xmin=0, xmax=21, ymin=0, ymax=7, xtick={0,5,10,15,20}, ytick={1,2,3,4,5,6}, xlabel={世代 $g$}, ylabel={$n_g/n_0$}]
  \addplot[bkcurve, domain=0:20, samples=21] {1.1^x};
  \addplot[bkthin, line width=1.4pt, dash pattern=on 4pt off 2pt, domain=0:20, samples=2] {1};
  \addplot[bkcurve2, dashed, domain=0:20, samples=21] {0.9^x};
  \node[font=\sffamily\footnotesize, text=bkBlue!80!black, anchor=south east] at (axis cs:19,6.2) {$k=1.1$（実線）};
  \node[font=\sffamily\footnotesize, text=bkGray, anchor=south] at (axis cs:15,1.1) {$k=1$（灰色の破線）};
  \node[font=\sffamily\footnotesize, text=bkOrange!80!black, anchor=north west] at (axis cs:12,0.95) {$k=0.9$（だいだいの破線）};
\end{axis}
\end{tikzpicture}
\bkfigcaption{世代ごとの中性子の数。$k=1.1$ では20世代で約6.7倍、$k=0.9$ では約0.12倍になる。}\label{ph:原-15:fig-gen}
\end{bkfigure}

\phboxhead{発見から原子炉まで}
核分裂が見つかるとすぐに、分裂で中性子が出れば連鎖反応が起こせることに、多くの物理学者が気づきました。1942年12月2日、アメリカのシカゴで、フェルミたちの組が、黒鉛（炭素）のブロックとウランを積み上げた装置で、はじめて連鎖反応を制御しながら続けることに成功しました（コラム）。一方、核分裂の発見から7年もたたない1945年には、連鎖反応を一気に起こす兵器（原子爆弾）が作られ、広島と長崎で使われました。このことは、原-16 で、放射線の健康への影響の調べ方とあわせて取り上げます。原子炉は、連鎖反応を $k=1$ に保ち、熱を少しずつ取り出し続けるための装置です。

\bksubsection{3}{原子炉のなかみ：燃料・減速材・冷却材・制御棒}
\label{ph:原-15:sec3}
いま世界でいちばん多く使われているのは、ふつうの水（軽水）を使う\textbf{軽水炉}です。\bkidx{けいすいろ}{軽水炉}日本の原子力発電所も軽水炉で、燃料はウラン235 の割合を3〜5\,\%に高めた（濃縮した）ウランです。原子炉は、次の4つでできています。

\begin{enumerate}[label={\bkcircled{\arabic*}}, leftmargin=2.4\zw, itemsep=1pt, topsep=2pt]
\item \textbf{燃料}：ウランを焼き固めた小さな円柱（ペレット）を、金属（ジルコニウムの合金）の細い管（被覆管）につめた燃料棒。燃料棒をたばにして並べる。分裂片は、ふつうはこの管の中にとどまる。
\item \textbf{減速材}：速い中性子をぶつけて遅くする物。軽水炉では水。水素の原子核（陽子）は中性子とほぼ同じ質量なので、1回の衝突で多くのエネルギーをうばう（原-13 ⑧）。\bkidx{げんそくざい}{減速材}
\item \textbf{冷却材}：燃料で生まれた熱を運び出す物。軽水炉では、減速材の水がそのまま冷却材になる。
\item \textbf{制御棒}：中性子をよく吸う物質（ホウ素など）でできた棒。燃料のあいだに出し入れして、吸われる中性子の数を変え、$k$ を調節する。\bkidx{せいぎょぼう}{制御棒}
\end{enumerate}

\begin{bkfigure}
\begin{tikzpicture}[x=1mm, y=1mm]
  % 圧力容器
  \fill[bkBlue!10] (6,4) rectangle (40,52);
  \draw[bkInk, line width=1.4pt] (6,4) rectangle (40,52);
  \foreach \x in {11,17,23,29,35} {\fill[bkOrange!60] (\x-1.3,8) rectangle (\x+1.3,32); \draw[bkInk, line width=0.4pt] (\x-1.3,8) rectangle (\x+1.3,32);}
  \foreach \x in {14,26} {\fill[bkGray!70] (\x-1,20) rectangle (\x+1,58); \draw[bkInk, line width=0.4pt] (\x-1,20) rectangle (\x+1,58);}
  \foreach \y in {36,40,44} {\draw[bkGray, line width=0.6pt] (9,\y) .. controls (12,\y+1.5) and (14,\y-1.5) .. (17,\y);}
  \node[txp lab, anchor=north] at (23,3) {圧力容器（中は水と蒸気）};
  \node[txp lab, anchor=south] at (20,58.5) {制御棒};
  % 蒸気の管とタービン
  \draw[bkRed!70!black, line width=2pt] (40,46) -- (66,46);
  \draw[bkRed!70!black, txp arr] (52,46) -- (58,46);
  \filldraw[fill=bkGray!25, draw=bkInk, line width=0.8pt] (66,40) -- (66,52) -- (82,56) -- (82,36) -- cycle;
  \node[txp lab] at (74,46) {タービン};
  \draw[bkInk, line width=1.2pt] (82,46) -- (88,46);
  \filldraw[fill=bkYellow!30, draw=bkInk, line width=0.8pt] (88,40) rectangle (104,52);
  \node[txp lab] at (96,46) {発電機};
  % 復水器
  \draw[bkRed!70!black, line width=1.6pt] (74,36) -- (74,28);
  \filldraw[fill=bkBlue!15, draw=bkInk, line width=0.8pt] (64,14) rectangle (86,28);
  \node[txp lab] at (75,21) {復水器};
  \draw[bkBlue!70!black, line width=1.4pt, -{Stealth[length=2.2mm]}] (100,17) -- (86,17);
  \draw[bkBlue!70!black, line width=1.4pt, -{Stealth[length=2.2mm]}] (86,25) -- (100,25);
  \node[txp lab, anchor=west] at (101,21) {海水};
  % もどる水
  \draw[bkBlue!80!black, line width=1.6pt] (64,18) -- (48,18) -- (48,12) -- (40,12);
  \draw[bkBlue!80!black, txp arr] (52,18) -- (50,18);
  \node[txp lab, anchor=north west] at (64,13) {ポンプで水をもどす};
\end{tikzpicture}
\bkfigcaption{沸騰水型の原子炉の発電のしくみ（形は簡単にしてある）。炉心で水を沸騰させ、その蒸気でタービンを回す。加圧水型では、炉心の水に高い圧力をかけて沸騰させず、その熱で別の水を蒸気にする（蒸気発生器）。どちらも、蒸気から先は火力発電所と同じランキンサイクル（熱-8）。}\label{ph:原-15:fig-plant}
\end{bkfigure}

熱でタービンを回す部分は、火力発電所とまったく同じです（熱-8）。原子力発電所の蒸気は約 280\,℃で、石炭火力の約 600\,℃より低いので、熱効率も約33\,\%と、火力より低くなります（熱-5 ⑦、この章の⑬）。中学理科で学ぶように、発電所は、水力・火力・原子力・太陽光など、いろいろなもとからエネルギーを取り出して電気に変えています。原子力の特徴は、熱のもとが燃焼ではなく核分裂であることです。

\phboxhead{なぜ減速材がいるのか}
2 MeV の中性子は、速さが光の約6.5\,\%（約 $\qty{2e7}{m/s}$）もあります。速い中性子は、ウラン235 の原子核のそばを、分裂を起こさずに通りすぎやすいのです。一方、熱中性子は約 \qty{2200}{m/s} とゆっくりで、分裂の断面積はとても大きくなります。そこで、燃料のあいだに水を置き、中性子を水素と十数回ぶつけて、熱中性子にしてから燃料にもどします。燃料と水の割合を変えると、中性子の遅くなり方が変わり、$k$ も変わります（⑤）。水は、減速材・冷却材・中性子のたて（遮へい）の3役をこなす、便利な物質なのです。

\Needspace{16\baselineskip}\phstage{te}{中性子を数えて、増倍率を作る}{3}

\bksubsection{4}{中性子の勘定：生まれる・吸われる・もれる}
\label{ph:原-15:sec4}
この本のはじめ（前付）の「保存則の地図」に、中性子の行がありました。「数える量：中性子、生じる量：生まれる数－吸収される数、法則：中性子のバランス、増倍率」です。ここで、その行を読み解きます。

原子炉の中の中性子の数 $n$ について、1秒あたりの勘定は、次のようになります。
\[ \text{たまる（1秒に増える数）}=\text{生まれる数}-\text{吸収される数}-\text{外へもれる数} \]
地図では、外へもれる数を、表の「入る・出るもの（中性子の流れ）」の列に入れ、「生じる量」には生まれる数と吸収される数の差だけを書きました。質量やエネルギーの勘定（力-13・熱-2）と同じ形で、違うのは、中性子は「生まれる」「消える」ことができる点です。\bkidx{ちゅうせいしのかんじょう}{中性子の勘定}

\begin{bkfigure}
\begin{tikzpicture}[x=1mm, y=1mm]
  \draw[bkTheme, line width=1.2pt, rounded corners=3mm, fill=bkTheme!5] (38,4) rectangle (90,36);
  \node[txp lab, align=center] at (64,25) {原子炉の中の\\中性子の数 $n$};
  \node[txp lab, align=center, text=bkGray] at (64,12) {1世代の時間 $\ell$\\（$10^{-4}$ s ほど）};
  \draw[bkGreen!60!black, line width=2pt, -{Stealth[length=3mm]}] (6,20) -- (36,20);
  \node[txp lab, anchor=south, align=center] at (19,21) {分裂で生まれる\\$k\,n/\ell$};
  \draw[bkRed!70!black, line width=2pt, -{Stealth[length=3mm]}] (92,26) -- (122,26);
  \node[txp lab, anchor=south] at (107,27) {吸収される};
  \draw[bkRed!70!black, line width=2pt, dashed, -{Stealth[length=3mm]}] (92,12) -- (122,12);
  \node[txp lab, anchor=north] at (107,11) {外へもれる};
  \draw[bkGray, line width=0.6pt] (124,26) -- (126,26) -- (126,12) -- (124,12);
  \node[txp lab, anchor=west, align=left] at (127,19) {合わせて\\$n/\ell$};
\end{tikzpicture}
\bkfigcaption{中性子の勘定。いまいる $n$ 個の中性子は、平均 $\ell$ 秒で、吸われるか外へもれていなくなる。そのあいだに、次の世代の中性子が $k$ 倍の $kn$ 個生まれる。}\label{ph:原-15:fig-balance}
\end{bkfigure}

中性子が生まれてから、吸われるか、もれるかしていなくなるまでの平均の時間を、\textbf{中性子の世代の時間} $\ell$ とします。いまいる $n$ 個は $\ell$ 秒で入れかわるので、1秒にいなくなる数は $n/\ell$ です。そのあいだに次の世代が $k$ 倍生まれるので、1秒に生まれる数は $kn/\ell$ です。勘定の式は
\pheq{\frac{\dd n}{\dd t}=\frac{k\,n}{\ell}-\frac{n}{\ell}=\frac{k-1}{\ell}\,n}{個/s ＝ 個/s}
となります。これは 6-4 の「増え方が今の量に比例する」式で、答えは
\pheq{n=n_0\,e^{(k-1)t/\ell}}{個 ＝ 個}
です。$k=1$ なら、$n$ は時間がたっても変わりません。これが\textbf{臨界}です。増倍率 $k$ が1より大きいか小さいかは、指数の肩の符号が正か負かと同じことです。数学 9-9 の言葉では、この式の極が $s=(k-1)/\ell$ にあり、臨界は「極がちょうど0にある状態」、制御棒は「極の場所を動かす道具」です。

\begin{phflow}[中性子の勘定から、連鎖反応の式へ]
\phflowrow{$\dfrac{\dd n}{\dd t}=\dfrac{k-1}{\ell}\,n$}{$k$ が一定\\1世代 $\ell$ ごとに見る}{$n_g=n_0\,k^g$}
\phflowrow{$\dfrac{\dd n}{\dd t}=\dfrac{k-1}{\ell}\,n$}{$k=1$\\（臨界）}{$n=$ 一定}
\end{phflow}

\phboxhead{表計算で、世代を追う}
表計算で、1世代ごとに $k$ 倍する表を作ってみましょう（表\ref{ph:原-15:tab-sheet}）。はじめを1000個とし、B列は $k=0.9$、C列は $k=1$、D列は $k=1.1$ です。

\begin{bkcenter}
\bksheet{4}{20}{%
1 & 世代 $g$ & $k=0.9$ & $k=1$ & $k=1.1$ \\
2 & 0 & 1000 & 1000 & 1000 \\
3 & 1 & 900 & 1000 & 1100 \\
4 & 2 & 810 & 1000 & 1210 \\
5 & 3 & 729 & 1000 & 1331 \\
6 & 4 & 656 & 1000 & 1464 \\
7 & 5 & 590 & 1000 & 1611 \\
8 & 10 & 349 & 1000 & 2594 \\}
\bktabcaption{世代ごとの中性子の数（小数は四捨五入）。B3 の式は \texttt{=B2*0.9}、D3 の式は \texttt{=D2*1.1} で、下へ貼りつける。8行めは10世代め（途中の行は省いた）。}\label{ph:原-15:tab-sheet}
\end{bkcenter}

1世代が $10^{-4}$ 秒なら、10世代は1000分の1秒です。$k=1.1$ では、1000分の1秒で2.6倍、1秒（1万世代）では、数えきれないほどの倍率になります。もし中性子が全部、分裂と同時に出るのなら、$k$ を0.1\,\%（$k=1.001$）まちがえただけで、出力は1秒で約2万倍になってしまいます（例題3）。人の手も機械も、とても間に合いません。それでも原子炉を一定に保てるのはなぜか。その答えのかぎが、⑥の\textbf{遅発中性子}です。

\bksubsection{5}{中性子の一生と4因子公式：なぜ濃縮が必要か}
\label{ph:原-15:sec5}
増倍率 $k$ の中身を、1個の中性子の「一生」をたどって調べます。燃料と水が十分に大きく広がっていて、外へもれる中性子がない原子炉を考え、その増倍率を $k_\infty$（ケイ・インフィニティ：無限に大きい炉の増倍率）と書きます。図\ref{ph:原-15:fig-life}は、速い中性子1000個が生まれてから、次の世代の中性子が生まれるまでを、順に追ったものです。

\begin{bkfigure}
\begin{tikzpicture}[x=1mm, y=1mm]
  \tikzset{lb/.style={draw=bkTheme, line width=0.9pt, fill=white, rounded corners=2mm, align=center, font=\sffamily\footnotesize, minimum height=11mm, text width=22mm, inner sep=1mm}}
  \node[lb] (a) at (13,40) {速い中性子\\\textbf{1000}};
  \node[lb] (b) at (47,40) {速い中性子\\\textbf{1030}};
  \node[lb] (c) at (81,40) {熱中性子\\\textbf{773}};
  \node[lb] (d) at (115,40) {燃料に吸われる\\\textbf{695}};
  \node[lb, draw=phFlowR, fill=phFlowRBg] (e) at (115,2) {次の世代\\\textbf{1282}};
  \draw[bk link] (a) -- node[above, font=\sffamily\footnotesize, text=bkInk, yshift=1mm] {$\times\varepsilon$} (b);
  \draw[bk link] (b) -- node[above, font=\sffamily\footnotesize, text=bkInk, yshift=1mm] {$\times p$} (c);
  \draw[bk link] (c) -- node[above, font=\sffamily\footnotesize, text=bkInk, yshift=1mm] {$\times f$} (d);
  \draw[bk link] (d) -- node[right, font=\sffamily\footnotesize, text=bkInk, xshift=1mm] {$\times\eta$} (e);
  \draw[bk link, dashed] (e.south) -- (115,-7) -- (13,-7) -- (a.south);
  \node[txp lab, anchor=north west, align=left, text width=84mm] at (20,30) {① $\varepsilon$：速いまま、ウラン238 を少し分裂させて増える。\\② $p$：遅くなる途中で、ウラン238 に吸われずに残る。\\③ $f$：熱中性子のうち、水などでなく燃料に吸われる。\\④ $\eta$：燃料に吸われた1個あたり、分裂で生まれる数。};
\end{tikzpicture}
\bkfigcaption{中性子の一生（ウラン235 が3\,\%の燃料。$\varepsilon=1.03$、$p=0.75$、$f=0.90$ は例として置いた値。$\eta=1.845$ は例題4 で計算する）。1000個から始めて、次の世代は1282個になる。}\label{ph:原-15:fig-life}
\end{bkfigure}

4つの段の割合をかけ合わせると、1世代での増え方がわかります。
\pheq{k_\infty=\eta\,\varepsilon\,p\,f}{1 ＝ 1 × 1 × 1 × 1}
\label{ph:原-15:4factor}これを\textbf{4因子公式}といいます。\bkidx{よんいんしこうしき}{4因子公式}4つの記号の意味は、次のとおりです。
\begin{itemize}[leftmargin=1.6\zw, itemsep=1pt, topsep=2pt]
\item $\eta$（イータ）：燃料に熱中性子が1個吸われたとき、分裂で生まれる中性子の平均の数。吸われても分裂せず、ただ捕まるだけのこともあるので、$\nu=2.43$ より小さい。
\item $\varepsilon$（イプシロン）：高速分裂因子。速い中性子が、ウラン238 を少し分裂させる分の増え方（1より少し大きい）。
\item $p$：共鳴を逃れる確率。中性子が遅くなる途中、ある決まったエネルギーのところで、ウラン238 にとても吸われやすくなる（共鳴吸収）。そこをくぐりぬけて熱中性子まで遅くなれる割合。
\item $f$：熱中性子利用率。熱中性子のうち、水や被覆管などでなく、燃料に吸われる割合。
\end{itemize}

\begin{phyomi}[$k_\infty=\eta\varepsilon pf$]
「次の世代の数は、関所を4つ通るごとに、それぞれの割合をかけたもの」と読みます。1000人が4つの関所を順に通り、それぞれで決まった割合だけが先へ進めるとき、最後に残る人数が、はじめの人数にすべての割合をかけた数になるのと同じです。どれか1つでも小さくなると、$k_\infty$ は小さくなります。実際の原子炉は大きさが有限なので、外へにげない割合（もれない確率）もかけ、$k=k_\infty\times$（もれない確率）とします（⑦）。
\end{phyomi}

\phboxhead{$\eta$ と濃縮度}
燃料が熱中性子を吸うとき、ウラン235 の分裂・ウラン235 の捕獲・ウラン238 の捕獲の3つが起こり得ます。ウラン235 とウラン238 の原子の数の割合を $N_5$・$N_8$ とすると、吸われる確率は断面積と数の積に比例するので、
\pheq{\eta=\nu\,\frac{N_5\,\sigma_{\mathrm f5}}{N_5(\sigma_{\mathrm f5}+\sigma_{\mathrm c5})+N_8\,\sigma_{\mathrm c8}}}{1 ＝ 1 × b/b}
です。$\sigma_{\mathrm f5}=\qty{585}{b}$（ウラン235 の分裂）、$\sigma_{\mathrm c5}=\qty{99}{b}$（ウラン235 の捕獲）、$\sigma_{\mathrm c8}=\qty{2.68}{b}$（ウラン238 の捕獲）を入れると、天然ウラン（ウラン235 が 0.72\,\%）では $\eta=1.35$、3\,\%に濃縮すると $\eta=1.85$ になります（例題4）。天然ウランのままでは、$\eta$ に $p$ と $f$（どちらも1より小さい）をかけると、ふつうの水を使う炉では $k_\infty$ が1にとどきません。そのため、軽水炉ではウラン235 の割合を3〜5\,\%に高める\textbf{濃縮}をします。\bkidx{のうしゅく}{濃縮}フェルミたちの最初の原子炉や、重水（水素の代わりに重水素をふくむ水）を使う原子炉は、中性子をほとんど吸わない減速材を使って $f$ を大きくし、天然ウランで臨界にしています。

\begin{phbaisuru}{濃縮度}
\noindent\begin{minipage}[c]{0.45\linewidth}\raggedright
濃縮度を2倍にしても、$\eta$ は2倍にはなりません。3\,\%を6\,\%にしても、$\eta$ は1.85から約1.95へ、5\,\%ほど増えるだけです。ウラン238 に吸われる分が小さくなると、残りはウラン235 自身の捕獲（$\sigma_{\mathrm c5}$）で決まり、どこまで濃縮しても $\eta$ は $2.43\times585/684=2.08$ をこえません。濃縮度を上げる効き目は、低い所ほど大きいのです。
\end{minipage}\hfill
\begin{minipage}[c]{0.52\linewidth}\centering
\begin{tikzpicture}
\begin{axis}[bkplot, width=56mm, height=40mm, xmin=0, xmax=21, ymin=0, ymax=2.4, xtick={0,5,10,15,20}, ytick={0.5,1,1.5,2}, xlabel={濃縮度（\%）}, ylabel={$\eta$}]
  \addplot[bkcurve, domain=0.3:20, samples=80] {2.43*585*x/(684*x+2.68*(100-x))};
  \addplot[bkthin, dashed, domain=0:20, samples=2] {2.078};
  \fill[bkBlue] (axis cs:0.72,1.35) circle (1.6pt) (axis cs:3,1.845) circle (1.6pt) (axis cs:6,1.95) circle (1.6pt);
  \node[font=\sffamily\footnotesize, anchor=west] at (axis cs:1.2,1.2) {天然 1.35};
  \node[font=\sffamily\footnotesize, anchor=south, text=bkGray] at (axis cs:14,2.1) {上限 2.08};
\end{axis}
\end{tikzpicture}
\end{minipage}
\end{phbaisuru}

\phboxhead{濃縮の段の数}
熱-14 の例題12 で、気体にしたウラン（六フッ化ウラン）を小さな穴に通すと、ウラン235 をふくむ軽い分子が、約1.0043倍の速さでもれ出すことを見ました。1回通すごとに、ウラン235 とウラン238 の数の比が1.0043倍になります。0.72\,\%を3\,\%にするには、この比を $(0.03/0.97)\div(0.0072/0.9927)\approx4.3$ 倍にしなければなりません。$1.0043^m=4.3$ を解くと、$m\approx340$ 段です（例題5）。穴を通す装置を何百段もつなぐのは大がかりなので、いまは、ガスを高速で回す遠心分離の装置を多くつないで濃縮するのが主な方法です。

\bksubsection{6}{反応度と遅発中性子：原子炉がゆっくり変わるわけ}
\label{ph:原-15:sec6}
$k$ が1からどれだけずれているかを、\textbf{反応度} $\rho$ で表します。\bkidx{はんのうど}{反応度}
\pheq{\rho=\frac{k-1}{k}}{1 ＝ 1}
$k$ が1に近いときは $\rho\approx k-1$ です。$k=1.001$ なら $\rho\approx0.001$ で、専門の本では、$10^{-5}$ を1として「100 pcm」と書くこともあります。

\phboxhead{遅れて出てくる中性子}
分裂で出る中性子のほとんどは、分裂の瞬間に出ます（\textbf{即発中性子}）。ところが、ごく一部の中性子は、分裂片の放射性崩壊のあとで、遅れて出てきます。たとえば、分裂片の臭素87 は半減期 \qty{55.7}{s} でβ崩壊し、できたクリプトン87 が、すぐに中性子を1個放り出します。このような中性子を\textbf{遅発中性子}、そのもとになる分裂片（臭素87 など）を\textbf{先行核}といいます。\bkidx{ちはつちゅうせいし}{遅発中性子}ウラン235 の分裂では、出る中性子の\textbf{0.65\,\%}（割合 $\beta=0.0065$）が遅発中性子です。先行核の半減期は0.2秒から約1分までいろいろで、平均すると、遅発中性子は分裂から約 \qty{12.7}{s} 遅れて出てきます。

\begin{bkfigure}
\begin{tikzpicture}[x=1mm, y=1mm]
  \draw[bkInk, line width=0.8pt, -{Stealth[length=2.4mm]}] (0,6) -- (124,6);
  \node[txp lab, anchor=north] at (118,5) {時間};
  \foreach \x/\l in {4/0,40/0.0001\,s,100/12.7\,s（平均）} {\draw[bkInk, line width=0.6pt] (\x,4.8) -- (\x,7.2); \node[txp lab, anchor=north] at (\x,4) {\l};}
  \txpnuc{(4,22)}{3.6}{bkGreen}
  \node[txp lab, anchor=south] at (4,26) {分裂};
  \draw[txp n] (8,23.5) -- (36,29); \draw[txp n] (8,22.5) -- (36,25);
  \node[txp lab, anchor=west, align=left] at (38,27) {即発中性子 99.35\,\%（すぐ次の世代へ）};
  \draw[bkOrange!80!black, line width=1pt, dashed] (7,19) -- (92,15);
  \txpnuc{(96,15)}{2.4}{bkOrange}
  \node[txp lab, anchor=north] at (52,15.5) {先行核（臭素87 など）が崩壊するまで待つ};
  \draw[txp n] (98.5,16) -- (108,22);
  \node[txp lab, anchor=south, align=center] at (110,22.5) {遅発中性子\\0.65\,\%};
\end{tikzpicture}
\bkfigcaption{即発中性子と遅発中性子（時間の目もりは、見やすくするため一様でない）。遅発中性子は数は少ないが、平均約13秒も遅れて出る。}\label{ph:原-15:fig-delayed}
\end{bkfigure}

\phboxhead{平均の世代の時間}
遅発中性子がほんの0.65\,\%でも、世代の時間の平均は大きく変わります。即発中性子の世代の時間を $\ell=\qty{1e-4}{s}$ とすると、100個の中性子のうち、99.35個は $10^{-4}$ 秒で、0.65個は約12.7秒で次の世代を作るので、平均は
\[ 0.9935\times0.0001+0.0065\times12.7\approx\qty{0.083}{s} \]
で、即発中性子だけのときの約800倍です。④の式の $\ell$ をこの0.083秒に置きかえると、出力が $e$ 倍（約2.7倍）になる時間、\textbf{炉の周期} $T$ は
\pheq{T\approx\frac{0.083\ \mathrm{s}}{\rho}}{s ＝ s ÷ 1}
となります。$\rho=0.001$ なら約83秒です。これなら、出力計を見ながら、制御棒をゆっくり動かして調節できます。⑧で、もっと正しい計算をします。

\begin{phyomi}[炉の周期 $T\approx0.083\,\mathrm{s}/\rho$]
「反応度が小さいほど、出力はゆっくり増える。遅発中性子が、1世代の時間を平均0.083秒まで引きのばしている」と読みます。ただし、これが使えるのは $\rho$ が $\beta=0.0065$ よりずっと小さいときだけです。$\rho\geqq\beta$ になると、即発中性子だけで $k\geqq1$ になってしまい（\textbf{即発臨界}）、世代の時間はふたたび $10^{-4}$ 秒の世界にもどります。そこで、反応度を $\beta$ で割った値を「ドル」とよび、原子炉は1ドル（$\rho=\beta$）よりずっと小さい反応度の範囲で運転するように作られています。
\end{phyomi}

\begin{phbaisuru}{反応度 $\rho$}
\noindent\begin{minipage}[c]{0.45\linewidth}\raggedright
この段の目安の式では、炉の周期 $T$ は $\rho$ に反比例するので、反応度を2倍にすると、出力が $e$ 倍になる時間は半分になります。$\rho=0.001$ で83秒、0.002 で41秒です。ただし、$\rho$ が $\beta$ に近づくと、周期は式よりもずっと速く短くなります（⑧の正しい式では28秒）。反応度は、小さく小さく入れるのが原則です。
\end{minipage}\hfill
\begin{minipage}[c]{0.52\linewidth}\centering
\begin{tikzpicture}
\begin{axis}[bkplot, width=56mm, height=42mm, xmin=0, xmax=0.0042, ymin=0, ymax=170, xtick={0.001,0.002,0.003,0.004}, xticklabels={0.001,0.002,0.003,0.004}, ytick={50,100,150}, xlabel={$\rho$}, ylabel={$T$（s）}]
  \addplot[bkcurve, domain=0.0005:0.004, samples=80] {0.083/x};
  \addplot[bkcurve2, dashed, domain=0.0005:0.004, samples=80] {(0.0065-x)/(0.08*x)};
  \fill[bkBlue] (axis cs:0.001,83) circle (1.6pt) (axis cs:0.002,41.5) circle (1.6pt);
  \node[font=\sffamily\footnotesize, text=bkBlue!80!black, anchor=west] at (axis cs:0.0011,95) {目安の式（実線）};
  \node[font=\sffamily\footnotesize, text=bkOrange!80!black, anchor=west] at (axis cs:0.0021,58) {⑧の式（破線）};
\end{axis}
\end{tikzpicture}
\end{minipage}
\end{phbaisuru}

\phkokomade

\Needspace{9\baselineskip}\phstage{shiki}{拡散方程式と動特性の式で読む}{5}

\bksubsection{7}{1群の拡散方程式と、臨界の大きさ}
\label{ph:原-15:sec7}
小さな原子炉では、中性子の多くが外へもれ、$k$ が1にとどきません。原子炉が臨界になるには、ある大きさ以上でなければならないのです。どれだけの大きさがいるかを、中性子の流れの式から求めます。

\phboxhead{中性子は、多い所から少ない所へ流れる}
中性子は、原子核にぶつかるたびに向きを変え、酔っぱらいの千鳥足のように動きます。こうした粒の集まりは、全体として、数の多い所から少ない所へ流れます。熱が温度の高い所から低い所へ流れる（熱-11 のフーリエの法則）のと同じで、1秒に1 m$^2$ を通る中性子の数 $J$ は、中性子束 $\phi$ の傾きに比例します。
\pheq{J=-D\,\frac{\dd\phi}{\dd x}}{1/(m$^2$·s) ＝ m × 1/(m$^3$·s)}
$D$ を\textbf{拡散係数}といいます（中性子の拡散係数は長さの次元をもつように決めます）。

\phboxhead{うすい板の中の勘定}
厚さ $\Delta x$、面積1 m$^2$ のうすい板の中の、1秒あたりの勘定（④）を書きます。外へもれる数は、出ていく流れと入ってくる流れの差 $J(x+\Delta x)-J(x)\approx\dfrac{\dd J}{\dd x}\Delta x=-D\,\dfrac{\dd^2\phi}{\dd x^2}\Delta x$ です（8-12 ④の「となりとの差」、熱-11 と同じ）。吸収される数は $\Sigma_{\mathrm a}\phi\,\Delta x$、分裂で生まれる数は $\nu\Sigma_{\mathrm f}\phi\,\Delta x$ です（$\Sigma_{\mathrm a}$・$\Sigma_{\mathrm f}$ は吸収と分裂の巨視的断面積：原-13 ⑧）。ちょうどつり合う形を探すために、生まれる数を $k$ でわって、つり合いの式を作ります。
\pheq{D\,\frac{\dd^2\phi}{\dd x^2}-\Sigma_{\mathrm a}\phi+\frac{\nu\Sigma_{\mathrm f}}{k}\,\phi=0}{1/(m$^3$·s)}
\label{ph:原-15:diffusion}これを\textbf{1群の拡散方程式}といいます。\bkidx{かくさんほうていしき}{拡散方程式}「1群」は、中性子のエネルギーを1つにまとめて考えたという意味です。$k$ は「生まれる数を何分の1にすれば、ちょうどつり合うか」を表すので、この $k$ が、この炉の増倍率です。

$k_\infty=\nu\Sigma_{\mathrm f}/\Sigma_{\mathrm a}$（もれのない炉の増倍率）、$L^2=D/\Sigma_{\mathrm a}$ を使って $\Sigma_{\mathrm a}$ でわると、
\pheq{\frac{\dd^2\phi}{\dd x^2}+B^2\phi=0,\qquad B^2=\frac{k_\infty/k-1}{L^2}}{1/m$^2$}
となります。弦の定在波（波-6）や箱の中の電子（原-5）と同じ形の式です。

\phboxhead{はだかの平板：$\cos$ の形}
厚さ $a$ の平らな炉心を考え、まん中を $x=0$ とします。外側には何もない（はだかの炉）ので、はしで中性子束は0とします。$\phi''=-B^2\phi$ の答えで、はしで0になり、どこでも正のものは
\[ \phi=\phi_0\cos\frac{\pi x}{a},\qquad B=\frac{\pi}{a} \]
です。弦のいちばん低い固有振動（波-6）と同じ半波長1つの形です。$B^2=(\pi/a)^2$ を、さきの $B^2$ の式に入れて $k$ について解くと、
\pheq{k=\frac{k_\infty}{1+L^2B^2}}{1 ＝ 1}
\label{ph:原-15:keff}が出ます。$1/(1+L^2B^2)$ が\textbf{もれない確率}で、$L^2B^2$ が、吸われる数に対するもれる数の割合です。臨界（$k=1$）の条件は
\pheq{B^2=\frac{k_\infty-1}{L^2}}{1/m$^2$}
です。$B^2$ を\textbf{バックリング}といいます。\bkidx{ばっくりんぐ}{バックリング}「中性子束の曲がりぐあい」のことで、形と大きさで決まります。速い中性子が遅くなるまでに進む距離もふくめるときは、$L^2$ のかわりに\textbf{移動面積} $M^2$ を使います。

\begin{phflow}[1群の拡散方程式から、臨界の大きさへ]
\phflowrow{$D\phi''-\Sigma_{\mathrm a}\phi+\dfrac{\nu\Sigma_{\mathrm f}}{k}\phi=0$}{はだかの平板\\はしで $\phi=0$}{$\phi=\phi_0\cos\dfrac{\pi x}{a}$}
\phflowrow{$k=\dfrac{k_\infty}{1+M^2B^2}$}{$k=1$（臨界）\\平板 $B=\pi/a$}{$a=\dfrac{\pi M}{\sqrt{k_\infty-1}}$}
\end{phflow}

\phboxhead{球と立方体}
球の炉心では、中性子束は中心からの距離 $r$ だけで決まります。8-12 ⑦の球のラプラシアンを使うと、拡散方程式は
\[ \frac{1}{r^2}\frac{\dd}{\dd r}\Bigl(r^2\frac{\dd\phi}{\dd r}\Bigr)+B^2\phi=0 \]
です。$\phi=u/r$ とおくと、$u''+B^2u=0$ になります（代入して確かめられます）。$u=\sin Br$ で、表面 $r=R$ で $\phi=0$ になるには $BR=\pi$ なので、
\[ \phi=\phi_0\,\frac{\sin(\pi r/R)}{\pi r/R},\qquad B=\frac{\pi}{R} \]
です。1辺 $a$ の立方体では、$x$・$y$・$z$ の3つの向きに $\cos$ の形が重なるので、$B^2=3(\pi/a)^2$ です。半径 $R$、高さ $H$ の円柱では、半径の向きにベッセル関数とよばれる関数が出てきて、
\[ B^2=\Bigl(\frac{2.405}{R}\Bigr)^2+\Bigl(\frac{\pi}{H}\Bigr)^2 \]
になります（2.405 は、波-6 ⑦の円い膜の振動で出た数と同じ）。

\begin{bkfigure}
\begin{tikzpicture}
\begin{axis}[bkplot, axis x line=bottom, axis y line=left, width=112mm, height=52mm, xmin=-1.15, xmax=1.15, ymin=0, ymax=1.2, xtick={-1,-0.5,0,0.5,1}, xticklabels={$-\frac{a}{2}$,,0,,$\frac{a}{2}$}, ytick={0.5,1}, xlabel={$x$}, ylabel={$\phi/\phi_0$}]
  \addplot[bkcurve, domain=-1:1, samples=80] {cos(deg(pi*x/2))};
  \addplot[bkcurve2, dashed, domain=0.005:1, samples=80] {sin(deg(pi*x))/(pi*x)};
  \addplot[bkcurve2, dashed, domain=-1:-0.005, samples=80] {sin(deg(pi*x))/(pi*x)};
  \node[font=\sffamily\footnotesize, text=bkBlue!80!black, anchor=south west] at (axis cs:0.12,1.12) {平板（実線）};
  \node[font=\sffamily\footnotesize, text=bkOrange!80!black, anchor=south east] at (axis cs:-0.12,1.12) {球（破線）};
\end{axis}
\end{tikzpicture}
\bkfigcaption{はだかの炉の中性子束の形（実線は平板の $\cos(\pi x/a)$、破線は球の $\sin(\pi r/R)/(\pi r/R)$）。はしでは中性子が外へにげるので少なく、まん中がいちばん多い。球は、横の軸を中心からの距離 $r$（$-R$ から $R$）と読む。球のほうが、まん中に集まった形になる。}\label{ph:原-15:fig-flux}
\end{bkfigure}

\begin{phmodel}[1群・はだかの炉]
\phignore{中性子のエネルギーによる断面積の違い（1つのエネルギーで代表させる）。燃料と水が分かれて並んでいること（全体を一様な物質とみなす）。炉心のまわりの反射体。}
\phassume{炉心は一様。中性子は拡散の式にしたがって動き、炉心のはしで中性子束は0。}
\phlimit{炉心のはしのすぐ近くや、制御棒のそば（拡散の式の精度が落ちる）。まわりに水などの反射体があるとき（中性子がはね返り、臨界の大きさは小さくてすむ）。}
\end{phmodel}

\begin{phbaisuru}{炉心の半径}
\noindent\begin{minipage}[c]{0.45\linewidth}\raggedright
球の炉心の半径 $R$ を2倍にすると、$B^2=(\pi/R)^2$ は $1/4$ になります。吸われる数に対するもれる数の割合 $M^2B^2$ も $1/4$ です。体積は8倍になるのに、中性子がにげる表面は4倍にしかならないからです（表面と体積の比）。$k_\infty=1.10$・$M^2=\qty{60}{cm^2}$ なら、半径 \qty{50}{cm} では $k=0.89$（未臨界）、\qty{100}{cm} では1.04 です。
\end{minipage}\hfill
\begin{minipage}[c]{0.52\linewidth}\centering
\begin{tikzpicture}
\begin{axis}[bkplot, width=56mm, height=42mm, xmin=0, xmax=210, ymin=0, ymax=0.3, xtick={50,100,150,200}, ytick={0.1,0.2}, yticklabels={0.1,0.2}, xlabel={$R$（cm）}, ylabel={$M^2B^2$}]
  \addplot[bkcurve, domain=45:200, samples=80] {60*(pi/x)^2};
  \fill[bkBlue] (axis cs:100,0.0592) circle (1.6pt) (axis cs:200,0.0148) circle (1.6pt);
  \node[font=\sffamily\footnotesize, anchor=south west] at (axis cs:100,0.065) {0.059};
  \node[font=\sffamily\footnotesize, anchor=south] at (axis cs:190,0.025) {0.015};
\end{axis}
\end{tikzpicture}
\end{minipage}
\end{phbaisuru}

\phboxhead{表計算・プログラムで確かめる}
平板を100のマスに区切り、8-12 ③・8-13 の「となりとの差」で拡散方程式を書くと、連立1次方程式になります。生まれる中性子を仮に置いて中性子束を求め、それから次の「生まれる中性子」を作る、というくり返しをすると、$\phi$ は $\cos$ の形に、$k$ は増倍率に近づいていきます。$D=\qty{3.0}{cm}$、$\Sigma_{\mathrm a}=\qty{0.05}{cm^{-1}}$（$L^2=\qty{60}{cm^2}$）、$\nu\Sigma_{\mathrm f}=\qty{0.055}{cm^{-1}}$（$k_\infty=1.10$）、厚さ $a=\qty{76.95}{cm}$（例題7 の臨界の厚さ）とします。
\begin{phcode}[Python]
import numpy as np
D, Sa, nSf = 3.0, 0.05, 0.055    \# 拡散係数 cm、Σa・νΣf（1/cm）
a, N = 76.95, 100                \# 平板の厚さ cm、区切りの数
h = a / N
M = np.zeros((N-1, N-1))         \# もれと吸収の行列（両はしで φ=0）
for i in range(N-1):
    M[i, i] = 2*D/h**2 + Sa
    if i > 0:   M[i, i-1] = -D/h**2
    if i < N-2: M[i, i+1] = -D/h**2
phi, k = np.ones(N-1), 1.0
for it in range(200):            \# 生まれる中性子から、次の φ を作る
    new = np.linalg.solve(M, nSf * phi / k)
    k = k * new.sum() / phi.sum()
    phi = new
print(round(k, 4))
print((phi / phi.max())[[9, 24, 49, 74, 89]].round(3))
\end{phcode}
\begin{phcodeout}[実行結果]
1.0
[0.309 0.707 1.    0.707 0.309]
\end{phcodeout}
$k$ はちょうど1になり（臨界）、中性子束は、はしから \qty{7.7}{cm}・\qty{19.2}{cm}・\qty{38.5}{cm}（まん中）… の所で 0.309・0.707・1 と、$\cos(\pi x/a)$ の値とそろいます。厚さを \qty{60}{cm} にすると $k=0.945$、\qty{100}{cm} にすると $k=1.039$ になり、どちらも $k_\infty/(1+M^2B^2)$ の式と小数第4位まで一致します。

\bksubsection{8}{1点炉の動特性：遅発中性子を式に入れる}
\label{ph:原-15:sec8}
⑥では遅発中性子の効き目を「平均の世代の時間」で見積もりました。ここでは、先行核の数を変数に加えて、出力の変わり方を正しく計算します。炉の中の中性子の数（出力に比例）を $n$、先行核の数を $C$ とします。先行核を1つのグループにまとめ、崩壊定数を $\lambda=\qty{0.08}{s^{-1}}$（平均の遅れ12.5秒ほど）とします。中性子の世代の時間を $\Lambda$（ラムダ）と書きます。

\begin{itemize}[leftmargin=1.6\zw, itemsep=1pt, topsep=2pt]
\item 1秒に生まれる中性子 $kn/\Lambda$ のうち、$(1-\beta)$ はすぐに出て、$\beta$ は先行核になる。
\item 先行核は、原-12 の崩壊の法則で、1秒に $\lambda C$ 個こわれ、そのたびに中性子が1個出る。
\end{itemize}
$k\approx1$ として、$k/\Lambda\approx1/\Lambda$、$(k-1)/\Lambda\approx\rho/\Lambda$ と近似すると、
\pheq{\frac{\dd n}{\dd t}=\frac{\rho-\beta}{\Lambda}\,n+\lambda C,\qquad \frac{\dd C}{\dd t}=\frac{\beta}{\Lambda}\,n-\lambda C}{1/s}
\label{ph:原-15:kinetics}となります。これを\textbf{1点炉の動特性の式}といいます。\bkidx{いってんろのどうとくせい}{1点炉の動特性}「1点」は、炉の中の場所による違いを考えず、炉全体を1つの点のように扱うという意味です。

\begin{phyomi}[1点炉の動特性の式]
1つめの式は「中性子は、即発の分で $(\rho-\beta)/\Lambda$ の割合で増え（$\rho<\beta$ なら実は減る）、先行核がこわれる分 $\lambda C$ で補われる」、2つめは「先行核は、分裂でできる分だけ増え、崩壊する分だけ減る」と読みます。2つの量がおたがいの増え方を決める、6-14 の「つながった2つの水そう」と同じ形です。運転中（$\rho=0$）は、2つめの式で $\dd C/\dd t=0$ から $C_0=\beta n_0/(\lambda\Lambda)$ で、先行核は中性子の約800倍もたまっています。
\end{phyomi}

\begin{phmodel}[1点炉・遅発中性子1群]
\phignore{炉の中の場所による出力の形の変化。先行核の半減期の違い（6つのグループを1つにまとめる）。温度による反応度の変化（⑩）。}
\phassume{反応度 $\rho$ は、ある時刻に階段のように変わり、その後は一定。$\Lambda=\qty{1e-4}{s}$、$\beta=0.0065$、$\lambda=\qty{0.08}{s^{-1}}$。}
\phlimit{大きな炉心で、場所によって出力の変わり方が違うとき。6つのグループを分けて計算すると、周期の値は2〜3割ずれる（$\rho=0.001$ で、1群では69秒、6群では53秒）。}
\end{phmodel}

\phboxhead{2つの指数の和になる}
6-14 ⑦の消去法で $C$ を消すと、$n$ だけの2階の微分方程式になります。答えは、2つの指数の和
\[ n=A_1e^{s_1t}+A_2e^{s_2t} \]
の形です。$s$ は、次の2次方程式の2つの解です。
\[ \Lambda s^2+(\beta-\rho+\lambda\Lambda)s-\lambda\rho=0 \]
$\Lambda$ がとても小さいので、2つの解は大きさがまったく違い、近似で
\pheq{s_1\approx\frac{\lambda\rho}{\beta-\rho},\qquad s_2\approx-\frac{\beta-\rho}{\Lambda}}{1/s}
です。$\rho=0.001$ では、$s_1=\qty{0.0145}{s^{-1}}$（周期 $T=1/s_1=\qty{69}{s}$）、$s_2=\qty{-55}{s^{-1}}$ です。$s_2$ の項は約0.02秒で消え、そのあいだに出力は、すぐに一段だけ跳び上がります。この跳びを\textbf{即発跳躍}といい、大きさは $\beta/(\beta-\rho)$（$\rho=0.001$ で1.18倍）です。\bkidx{そくはつちょうやく}{即発跳躍}そのあとは、$s_1$ の項がゆっくりと増えていきます。9-9 の言葉では、$s_1$ と $s_2$ がこの炉の2つの極で、反応度が0のときは $s_1=0$（極がちょうど0）、反応度を負にすると $s_1$ が左（負）へ動いて、出力は減っていきます。

\phboxhead{ルンゲ・クッタ法で確かめる}
2つの式を、6-2 のもっと知りたいのルンゲ・クッタ法（4次）で、同じ時刻の値を使って同時に進めます（6-14 ③）。
\begin{phcode}[Python]
import numpy as np
rho, beta, lam, Lam = 0.001, 0.0065, 0.08, 1e-4
def f(y):                        \# y = [n, C]
    n, C = y
    return np.array([(rho - beta) / Lam * n + lam * C,
                     beta / Lam * n - lam * C])
y = np.array([1.0, beta / (Lam * lam)])   \# はじめはつり合い（n=1）
h = 1e-4                         \# 1歩の時間（s）
for i in range(1, 600001):       \# 60 秒まで
    k1 = f(y); k2 = f(y + h/2*k1)
    k3 = f(y + h/2*k2); k4 = f(y + h*k3)
    y = y + h/6*(k1 + 2*k2 + 2*k3 + k4)
    if i in (1000, 10000, 100000, 600000):
        print(round(i*h, 1), round(y[0], 3))
\end{phcode}
\begin{phcodeout}[実行結果]
0.1 1.182
1.0 1.198
10.0 1.366
60.0 2.823
\end{phcodeout}
0.1秒で1.18倍に跳び（即発跳躍）、そのあとは69秒で $e$ 倍のペースでゆっくり増え、60秒で2.8倍です。1歩の時間 $h$ を $\Lambda$ と同じくらいに小さくしないと、$s_2$ の速い変化で計算がこわれてしまう（6-2 ⑥）ことにも注意しましょう。

\begin{bkfigure}
\begin{tikzpicture}
\begin{axis}[bkplot, width=112mm, height=50mm, xmin=0, xmax=63, ymin=0, ymax=3.2, xtick={10,20,30,40,50,60}, ytick={1,2,3}, xlabel={$t$（s）}, ylabel={$n/n_0$}]
  \addplot[bkcurve, domain=0:60, samples=200] {1.181*exp(0.01452*x)-0.1812*exp(-55.09*x)};
  \addplot[bkthin, dashed, domain=0:60, samples=2] {1};
  \addplot[bkcurve2, dashed, domain=0:0.116, samples=20] {exp(10*x)};
  \node[font=\sffamily\footnotesize, text=bkBlue!80!black, anchor=north west] at (axis cs:30,1.75) {遅発中性子がある（実線）};
  \node[font=\sffamily\footnotesize, text=bkOrange!80!black, anchor=west, align=left] at (axis cs:1.5,2.9) {即発中性子だけなら、0.1 秒で $e$ 倍、\\0.12 秒で 3.2 倍をこえる（破線）};
  \node[font=\sffamily\footnotesize, anchor=north west] at (axis cs:3,0.9) {即発跳躍で 1.18 倍（$t=0$ のすぐあと）};
\end{axis}
\end{tikzpicture}
\bkfigcaption{反応度 $\rho=0.001$ を階段のように入れたときの出力（1点炉、遅発中性子1群）。実線は $1.181e^{0.0145t}-0.181e^{-55t}$。}\label{ph:原-15:fig-kinetics}
\end{bkfigure}

\bksubsection{9}{はてな？の答え：遅れる中性子と、自分で下がる増倍率}
\label{ph:原-15:sec9}
\phboxhead{\bkhatenaref の答え}
原子炉は、次の3つのしくみで、反応を一定に保っています。
\begin{enumerate}[label={\bkcircled{\arabic*}}, leftmargin=2.4\zw, itemsep=1pt, topsep=2pt]
\item \textbf{制御棒で $k$ を1に合わせる}：ウラン235 が1回分裂すると平均2.43個の中性子が出ます。そのうち、ちょうど1個が次の分裂を起こすように、制御棒（中性子を吸う物）の出し入れで、吸われる中性子の数を調節します。$k=1$ が臨界です（②④）。
\item \textbf{遅発中性子が、変化をゆっくりにする}：1世代は $10^{-4}$ 秒ほどですが、中性子の約0.65\,\%は、分裂から平均約13秒遅れて出ます。そのため、$k=1.001$（反応度0.001）でも、出力が $e$ 倍になるのに約1分かかります。人や機械が出力計を見ながら制御棒を動かす時間が、十分にあります（⑥⑧）。
\item \textbf{温度が上がると、$k$ が自然に下がる}：出力が増えて燃料や水の温度が上がると、ウラン238 が中性子を吸いやすくなり、水の減速の働きも弱くなるので、$k$ がひとりでに下がるように作ってあります（負の反応度係数：⑩）。出力が増えすぎると、ブレーキがかかるのです。
\end{enumerate}

\begin{bkfigure}
\begin{tikzpicture}[x=1mm, y=1mm]
  \node[bk node, minimum width=24mm, minimum height=10mm, font=\sffamily\footnotesize] (r) at (14,22) {制御棒};
  \node[draw=bkInk, circle, line width=0.8pt, minimum size=7mm, font=\sffamily\footnotesize, inner sep=0pt] (s) at (40,22) {$+$};
  \node[bk node, minimum width=28mm, minimum height=10mm, font=\sffamily\footnotesize, align=center] (k) at (66,22) {反応度 $\rho$\\（$k$）};
  \node[bk node, minimum width=26mm, minimum height=10mm, font=\sffamily\footnotesize] (p) at (106,22) {出力};
  \node[bk node, minimum width=36mm, minimum height=10mm, font=\sffamily\footnotesize, align=center] (t) at (106,2) {燃料・水の温度が上がる};
  \draw[bk link] (r) -- (s); \draw[bk link] (s) -- (k); \draw[bk link] (k) -- (p);
  \draw[bk link] (p) -- (t);
  \draw[bk link] (t.west) -| node[pos=0.8, left, font=\sffamily\footnotesize, text=bkInk] {$-$} (s.south);
  \node[txp lab, anchor=south] at (64,3) {負のフィードバック};
\end{tikzpicture}
\bkfigcaption{原子炉の負のフィードバック（波-4 のブロック線図と同じ書き方）。出力が上がると温度が上がり、反応度を下げる向きにはたらく。}\label{ph:原-15:fig-feedback}
\end{bkfigure}

この3つのうち、②は物理の性質そのもの、③は燃料と水の配置の設計で決まります。エアコンが温度計を見て冷やし方を変える（波-4）ように、原子炉は、自分の温度を見て増倍率を変える、負のフィードバックを、はじめから物理として持っているのです。

\bksubsection{10}{反応度係数とキセノン：運転中に $k$ を変えるもの}
\label{ph:原-15:sec10}
\phboxhead{温度で変わる反応度}
温度が1 K 上がったときの反応度の変化を\textbf{反応度係数}（温度係数）といいます。\bkidx{はんのうどけいすう}{反応度係数}主なものは2つです。
\begin{itemize}[leftmargin=1.6\zw, itemsep=1pt, topsep=2pt]
\item \textbf{ドップラー係数}（燃料の温度）：ウラン238 が中性子を強く吸う「共鳴」のエネルギーは、とてもせまい範囲にあります。燃料の温度が上がると、ウラン238 の原子核の熱運動が激しくなり、中性子から見た共鳴のはばが広がって（音のドップラー効果と同じ考え：波-8）、吸われる中性子が増えます。$p$ が小さくなり、$k$ が下がります。燃料の温度が上がるとすぐに効くので、いちばん頼りになるブレーキです。
\item \textbf{減速材の温度係数}：水の温度が上がると水がふくらみ、沸騰すると泡（ボイド）ができて、減速する水素が減ります。軽水炉は、水が減ると中性子が遅くなりにくくなって $k$ が下がるように、燃料と水の割合を決めてあります。
\end{itemize}
温度係数が $-\qty{3}{pcm/K}$（$-3\times10^{-5}$ /K）なら、反応度0.001 を入れても、温度が約 \qty{33}{K} 上がれば打ち消されます（例題10）。

\phboxhead{キセノン：止めたあとにたまる「中性子の毒」}
分裂片の中には、中性子をとても強く吸うものがあります。いちばん大きいのがキセノン135 で、熱中性子の吸収断面積は約 $\qty{2.6e6}{b}$、ウラン235 の分裂の4000倍以上です。キセノン135 は、分裂で直接できる分（約0.2\,\%）より、分裂片のヨウ素135（半減期 \qty{6.58}{h}）がβ崩壊してできる分（約6\,\%）のほうがずっと多く、キセノン135 自身も半減期 \qty{9.14}{h} でこわれます。原-12 ⑧の「親と子」の形です。
\[ \frac{\dd I}{\dd t}=\gamma_I\Sigma_{\mathrm f}\phi-\lambda_II,\qquad \frac{\dd X}{\dd t}=\gamma_X\Sigma_{\mathrm f}\phi+\lambda_II-\lambda_XX-\sigma_X\phi X \]
$I$・$X$ は1 cm$^3$ あたりのヨウ素135・キセノン135 の数、$\gamma_I$・$\gamma_X$ は分裂1回あたりにできる割合、$\sigma_X\phi X$ は中性子を吸って消える数です。運転中は、キセノン135 の大部分（中性子束 $\qty{3e13}{cm^{-2}.s^{-1}}$ とすると約8割）が、崩壊する前に中性子を吸って消えています。

\noindent\begin{minipage}[c]{0.48\linewidth}\raggedright
原子炉を止めると（$\phi=0$）、中性子を吸って消える道がなくなるのに、ヨウ素135 からは、まだキセノン135 が作られ続けます。だから、止めてからしばらく、キセノン135 は運転中より\textbf{増えます}。$\phi=0$ では式は原-12 ⑧の親と子そのもので、例の値では、約8時間後に運転中の約1.9倍の山になり、約26時間後にもとの量にもどります（図\ref{ph:原-15:fig-xenon}・例題11）。この間は、中性子がたくさん吸われるので、原子炉をもう一度臨界にするのがむずかしくなります。
\end{minipage}\hfill
\begin{minipage}[c]{0.49\linewidth}\centering
\begin{tikzpicture}
\begin{axis}[bkplot, width=50mm, height=44mm, xmin=0, xmax=52, ymin=0, ymax=2.2, xtick={10,20,30,40,50}, ytick={0.5,1,1.5,2}, xlabel={$t$（h）}, ylabel={$X/X_0$}]
  \addplot[bkcurve, domain=0:48, samples=120] {exp(-0.07584*x)+3.263*0.10534/(0.07584-0.10534)*(exp(-0.10534*x)-exp(-0.07584*x))};
  \addplot[bkthin, dashed, domain=0:48, samples=2] {1};
  \fill[bkBlue] (axis cs:8.35,1.881) circle (1.6pt);
  \node[font=\sffamily\footnotesize, anchor=south west] at (axis cs:9,1.88) {8.3 h で 1.88 倍};
\end{axis}
\end{tikzpicture}
\bkfigcaption{止めたあとのキセノン135（運転中を1とした。例題11 の値）。}\label{ph:原-15:fig-xenon}
\end{minipage}

\bksubsection{11}{崩壊熱：止めたあとも、熱は出続ける}
\label{ph:原-15:sec11}
制御棒を全部入れて連鎖反応を止めても、燃料の中には、運転中にできた分裂片がたくさん残っています。分裂片は放射性で、β崩壊やγ線を出しながら、何秒・何日・何年もかけてこわれ続けます。そのエネルギーが熱になったものを\textbf{崩壊熱}といいます。\bkidx{ほうかいねつ}{崩壊熱}表\ref{ph:原-15:tab-energy}の「あとから出る」β と γ は、熱になる分の約7\,\%でした。運転中の原子炉の熱の約7\,\%は、もともと崩壊熱だったのです。だから、止めた直後にも、運転中の6〜7\,\%の熱が出ています。

崩壊熱の割合 $P/P_0$（$P_0$ は運転中の熱出力）は、何百種類もの分裂片の崩壊を足し合わせて計算します。目安には、ウェイとウィグナーが1948年に示した形の式がよく使われます。
\pheq{\frac{P}{P_0}\approx0.066\left[\,t^{-0.2}-(t+T_0)^{-0.2}\,\right]}{単位なし、$t$ は秒}
\label{ph:原-15:decayheat}$t$ は止めてからの時間（秒）、$T_0$ は運転した時間（秒）です。係数は本によって 0.062〜0.066 と少し違い、精密な値とは2倍ほどずれることもある、あくまで目安の式です。

\begin{phyomi}[{$P/P_0\approx0.066\,[t^{-0.2}-(t+T_0)^{-0.2}]$}]
「崩壊熱は、止めてからの時間の $-0.2$ 乗で、ゆっくり減っていく」と読みます。半減期の短い分裂片が先にこわれ、長い分裂片が残るので、1つの指数関数（原-12）ではなく、もっとゆっくりした減り方になります。1年運転した 3 GW の原子炉なら、1秒後に約 \qty{190}{MW}（6.4\,\%）、1時間後に約 \qty{32}{MW}（1.1\,\%）、1日後でも約 \qty{14}{MW}（0.47\,\%）です（例題12）。\qty{14}{MW} は、100\,℃の水を1時間に約22トン蒸発させる熱です。
\end{phyomi}

\begin{bkfigure}
\begin{tikzpicture}
\begin{axis}[bkplot, xmode=log, ymode=log, width=112mm, height=56mm, xmin=0.8, xmax=5e7, ymin=0.01, ymax=10, xtick={1,60,3600,86400,2592000,31557600}, xticklabels={1秒,1分,1時間,1日,30日,1年}, ytick={0.01,0.1,1,10}, yticklabels={0.01,0.1,1,10}, xlabel={止めてからの時間}, ylabel={$P/P_0$（\%）}, axis lines=left, grid=both, minor grid style={draw=none}]
  \addplot[bkcurve, domain=1:3.2e7, samples=120] {6.6*(x^(-0.2)-(x+31557600)^(-0.2))};
  \fill[bkBlue] (axis cs:1,6.39) circle (1.8pt);
  \fill[bkBlue] (axis cs:3600,1.074) circle (1.8pt);
  \fill[bkBlue] (axis cs:86400,0.471) circle (1.8pt);
  \fill[bkBlue] (axis cs:31557600,0.0270) circle (1.8pt);
  \node[font=\sffamily\footnotesize, anchor=south west] at (axis cs:1.3,6.4) {6.4 \%};
  \node[font=\sffamily\footnotesize, anchor=south west] at (axis cs:4000,1.1) {1.1 \%};
  \node[font=\sffamily\footnotesize, anchor=south west] at (axis cs:95000,0.48) {0.47 \%};
  \node[font=\sffamily\footnotesize, anchor=north east] at (axis cs:2.2e7,0.026) {0.027 \%};
\end{axis}
\end{tikzpicture}
\bkfigcaption{崩壊熱の割合（1年運転したあと。ウェイ・ウィグナーの形の目安の式）。たて・横とも、目もりが10倍ずつの対数の目もり（3-12）。}\label{ph:原-15:fig-decayheat}
\end{bkfigure}

\begin{phbaisuru}[10倍]{止めてからの時間}
止めてからの時間を10倍にしても、崩壊熱は $10^{-0.2}=0.63$ 倍、つまり約6割にしかなりません（運転した時間よりずっと短いあいだ）。1時間後の \qty{32}{MW} は、10時間後でも約 \qty{18}{MW} です。半減期のように「決まった時間で半分」にはならず、減り方がだんだんゆっくりになるのが、崩壊熱の特徴です。止めたあとも、何日も、何年も、冷やし続けなければならないのは、このためです。使い終わった燃料をプールで何年も冷やす理由も、ここにあります（原-16）。
\end{phbaisuru}

\phboxhead{福島第一原子力発電所の事故の物理}
2011年3月11日、東北地方太平洋沖地震が起きたとき、運転中だった東京電力福島第一原子力発電所の1〜3号機は、すべて自動で停止しました。連鎖反応は止まりました。しかし、崩壊熱は出続けます。地震で送電の鉄塔が倒れるなどして外からの電気が止まり、非常用のディーゼル発電機が動き出しましたが、そのあとの津波で、発電機や配電盤が水につかり、6号機を除いて交流の電源がすべて失われました。原子炉を冷やす設備が止まり、炉心に水を入れられない時間が、1号機で約14時間、2号機・3号機で約6時間半続いたとみられています。

水が減って燃料が水から出ると、崩壊熱で燃料は高い温度になり、溶けました（炉心溶融）。高い温度では、燃料の被覆管の金属ジルコニウムが水蒸気と反応して（\mbox{Zr＋2H$_2$O → ZrO$_2$＋2H$_2$}）、水素ができます。この水素が原子炉の建屋にたまり、1号機・3号機と、3号機から水素が流れこんだとみられる4号機で、水素爆発が起きました。放射性物質が外へ出て、多くの人が避難しました。事故の大きさは、国際原子力事象評価尺度で、いちばん深刻なレベル7とされています。

止めた原子炉でも、崩壊熱は、1日後で運転中の0.5\,\%ほど、熱出力の大きな原子炉では何 MW にもなります。冷やす手段を長い時間失うことが、どれほど大きな結果につながるかを、この事故は示しました。事故の経過と、廃炉の作業、地域の暮らしの今については、原-16 で取り上げます。

\bksubsection{12}{核融合の入口：軽い原子核をくっつける}
\label{ph:原-15:sec12}
図\ref{ph:原-15:fig-BA}の曲線で、鉄より軽い原子核は、くっつくと1核子あたりの結合エネルギーが大きくなります。重水素（D）と三重水素（T）がくっついてヘリウム4 と中性子になる反応
\[ \mathrm{D}+\mathrm{T}\ \to\ {}^4\mathrm{He}\,(\qty{3.5}{MeV})+\mathrm{n}\,(\qty{14.1}{MeV}) \]
では、1回に \qty{17.6}{MeV} が出ます（原-11 例題3）。運動量の保存から、軽い中性子がエネルギーの約8割を持っていきます。燃料1 g あたりでは、ウラン235 の分裂の約4倍のエネルギーです（例題13）。

ところが、DとTはどちらも正の電気を持つので、近づくとクーロン力で強く反発します。核力が届く距離まで近づけるには、燃料を1億度をこえる温度のプラズマ（原子核と電子がばらばらの気体）にして、激しくぶつけ合わせなければなりません。しかも、その熱いプラズマを、十分な密度で、十分に長く閉じこめる必要があります。点火（核融合の熱だけで温度を保てる状態）には、およそ
\[ n\,\tau_E\geqq\qty{1.6e20}{m^{-3}.s}\qquad（\text{温度は約 }\qty{26}{keV}\text{、約3億 K のとき）} \]
が必要と見積もられています（\textbf{ローソンの条件}の1つの形。$n$ は燃料の数密度、$\tau_E$ はエネルギーが逃げるまでの時間）。\bkidx{ろーそんのじょうけん}{ローソンの条件}

核融合は、核分裂のような連鎖反応ではありません。温める力やプラズマの閉じこめが少しでもくずれると、反応はすぐに止まります。一方、反応で出る \qty{14}{MeV} の中性子は、炉のかべの材料を傷め、材料を放射化します（原-13 ⑨）。日本もふくむ国々が南フランスで建設している実験炉 ITER は、燃料を温めるために入れる \qty{50}{MW} の10倍の、\qty{500}{MW} の核融合の出力（$Q=10$）をめざしています（熱-16）。

\begin{phmada}[核融合で、電気をつくれるか？]
2022年12月、アメリカの国立点火施設（NIF）で、レーザーで燃料の粒を一気に押しつぶす方式により、燃料に届けたレーザーのエネルギー（\qty{2.05}{MJ}）より多い核融合のエネルギー（\qty{3.15}{MJ}）が、はじめて出ました。ただし、そのレーザーを動かすのに使った電気のエネルギーは、それよりずっと大きいので、施設全体でエネルギーが増えたわけではありません。磁場でプラズマを閉じこめる方式では、ITER が $Q=10$ をめざして建設中です。入れたエネルギーより多くの\textbf{電気}を、長い期間、安定して取り出せる核融合の発電所は、まだできていません。プラズマを長く安定に閉じこめる方法、\qty{14}{MeV} の中性子に長く耐える材料、燃料の三重水素を炉の中で作り続ける方法など、いくつもの問題の研究が続いています。
\end{phmada}

\bksubsection{13}{原子力発電の熱の流れ}
\label{ph:原-15:sec13}
熱出力 \qty{3}{GW} の原子炉で、電気出力は約 \qty{1}{GW} です（熱効率約33\,\%：熱-5 例題）。残りの約 \qty{2}{GW} の熱は、復水器で海水に捨てられます。約280\,℃の蒸気と約30\,℃の海水の間のカルノーの効率は約45\,\%で（熱-5 ⑦）、実際の効率はその7割ほどです。

\begin{bkfigure}
\begin{tikzpicture}[x=1mm, y=1mm]
  \fill[bkRed!35] (0,10) rectangle (34,40);
  \node[txp lab, align=center] at (17,25) {核分裂の熱\\\qty{3.0}{GW}\\（100\,\%）};
  \fill[bkRed!35] (34,10) -- (34,40) -- (60,40) -- (60,10) -- cycle;
  \fill[bkYellow!60] (60,30) rectangle (120,40);
  \fill[bkBlue!25] (60,10) rectangle (120,30);
  \draw[bkInk, line width=0.5pt] (60,30) -- (120,30);
  \node[txp lab, anchor=west] at (62,35) {電気 約 \qty{1.0}{GW}（33\,\%）};
  \node[txp lab, anchor=west, align=left] at (62,20) {復水器から海へ\\約 \qty{2.0}{GW}（67\,\%）};
  \node[txp lab, anchor=south] at (47,41) {タービン};
  \draw[bkGray, line width=0.6pt, dashed] (60,6) -- (60,44);
\end{tikzpicture}
\bkfigcaption{熱出力 \qty{3}{GW} の原子力発電所の熱の流れ（目安）。海水の温度の上がり方 7 K は、例として置いた値。}\label{ph:原-15:fig-heatflow}
\end{bkfigure}

海へ捨てる熱 \qty{2.0}{GW} を、海水の温度を 7 K 上げて運ぶとすると、1秒あたりの海水の量は $2.0\times10^9\div(4.18\times10^3\times7)\approx\qty{6.8e4}{kg}$、約68トンです（水の比熱 \qty{4.18}{kJ/(kg.K)}：熱-1）。発電所の効率が、海の環境にもつながっていることがわかります。燃料の量で見ると、\qty{3}{GW} の原子炉が1日に分裂させるウラン235 は約 \qty{3}{kg} です（例題2）。

\begin{senmonsho}[原子炉物理の専門書の書き方]
\begin{bkhonbun}
一群拡散方程式 $D\nabla^2\phi-\Sigma_a\phi+\dfrac{1}{k_{\mathrm{eff}}}\nu\Sigma_f\phi=0$ において，裸の炉の幾何学的バックリング $B_g^2$ と材料バックリング $B_m^2=(\nu\Sigma_f-\Sigma_a)/D$ が等しいとき，炉は臨界である．実効増倍率は $k_{\mathrm{eff}}=k_\infty/(1+M^2B_g^2)$ で与えられる．また，一点炉動特性方程式において，即発臨界 $\rho=\beta$ を超えると炉周期は即発中性子寿命で支配される．
\end{bkhonbun}
\bkyaku{「裸の炉」は、まわりに反射体のない炉（この章のはだかの炉）。$k_{\mathrm{eff}}$（ケイ・エフ）は実効増倍率で、この章の $k$。$B_g^2$ は形と大きさで決まるバックリング、$B_m^2$ は材料で決まるバックリングで、その2つが等しいことが、この章の「臨界の条件」。$\nabla^2$ は、平板なら $\dd^2/\dd x^2$。「炉周期」はこの章の炉の周期 $T$、「即発中性子寿命」はこの章の $\ell$。}
\end{senmonsho}

\begin{phsaki}{原子炉の物理の先へ}
この章では、中性子のエネルギーを1つにまとめ（1群）、遅発中性子も1つにまとめて考えました。実際の原子炉の設計では、中性子のエネルギーを何十〜何百のグループに分けた\textbf{多群の拡散方程式}や、中性子の向きまで考える\textbf{輸送方程式}を、コンピュータで解きます。乱数で中性子を1個ずつ追いかけるモンテカルロ法（3-17）もよく使われます。遅発中性子は6つ以上のグループに分け、燃料の燃え方（燃焼）、熱と水の流れ（熱水力）、事故のときのふるまいまで、組み合わせて計算します。これらは、原子炉物理・原子炉工学・核融合工学の専門の本で学びます。
\end{phsaki}

\begin{taisetsu}
\def\tsl#1{\makebox[7.2\zw][l]{\sffamily\bfseries\small\color{\bkTT}#1}}
\tsl{核分裂}1回に約 \qty{200}{MeV}、中性子 $\nu=2.43$ 個（ウラン235・熱中性子）\\[3pt]
\tsl{連鎖反応}$n_g=n_0k^g$、$\dfrac{\dd n}{\dd t}=\dfrac{k-1}{\ell}n$、臨界は $k=1$\\[3pt]
\tsl{4因子公式}$k_\infty=\eta\varepsilon pf$、$k=\dfrac{k_\infty}{1+M^2B^2}$\\[3pt]
\tsl{反応度}$\rho=\dfrac{k-1}{k}$、遅発中性子 $\beta=0.0065$、周期 $T\approx\dfrac{\beta-\rho}{\lambda\rho}$\\[3pt]
\tsl{拡散方程式}$D\phi''-\Sigma_{\mathrm a}\phi+\dfrac{\nu\Sigma_{\mathrm f}}{k}\phi=0$、臨界 $B^2=\dfrac{k_\infty-1}{M^2}$\\[3pt]
\tsl{崩壊熱}$\dfrac{P}{P_0}\approx0.066\,[t^{-0.2}-(t+T_0)^{-0.2}]$（目安）
\end{taisetsu}

\begin{kotoba}
\small
\bkword{核分裂}{重い原子核が、ほぼ半分の大きさの2つの原子核に分かれること。}
\bkword{連鎖反応}{分裂で出た中性子が次の分裂を起こし、反応がつながっていくこと。}
\bkword{増倍率}{ある世代の中性子の数に対する、次の世代の数の比 $k$。$k=1$ が臨界。}
\bkword{減速材}{中性子をぶつけて遅くする物質（軽水炉では水）。}
\bkword{遅発中性子}{分裂片の放射性崩壊のあとで、遅れて出てくる中性子。ウラン235 では約0.65\,\%。}
\bkword{反応度}{$k$ の1からのずれ $\rho=(k-1)/k$。$\beta$ でわった値をドルとよぶ。}
\bkword{崩壊熱}{止めた原子炉の中で、分裂片の放射性崩壊が出し続ける熱。}
\end{kotoba}

\bkkigocards{\bkkigocard{$k$}{ケイ}{増倍率。次の世代と今の世代の中性子の数の比（原-15）。$k_\infty$ はもれのない炉の値。波数・ばね定数・クーロンの定数の $k$ とは別}{multiplication factor}{原-15}}
{\bkkigocard{$\rho$}{ロー}{反応度 $(k-1)/k$。単位なし。原-15 では密度を記号にしない}{reactivity}{原-15}}
\bkkigocards{\bkkigocard{$\nu$}{ニュー}{1回の分裂で出る中性子の平均の数（ウラン235 で2.43）。高校の教科書の光の振動数 $\nu$ は、この本では $f$}{neutrons per fission}{原-15}}
{\bkkigocard{$\beta$}{ベータ}{遅発中性子の割合（ウラン235 で0.0065）。原-11 の $v/c$、β線とは別}{delayed neutron fraction}{原-15}}
\bkkigocards{\bkkigocard{$B^2$}{ビー・二乗}{バックリング〔1/m$^2$〕。中性子束の曲がりぐあい。原-11 の結合エネルギー $B$ とは別}{buckling}{原-15}}{}

%% ======================== 例題 ========================
\Needspace{10\baselineskip}\bksection{reidai}{例題でたしかめよう}
\label{ph:原-15:ex1}
\begin{phreidai}[ウラン235 の 1 g が全部分裂したら]
ウラン235 の \qty{1}{g} が全部分裂すると、何 J のエネルギーが出ますか。それは、\qty{1}{MW} の熱を何日出し続けられる量ですか。また、炭素（燃焼熱 \qty{32.8}{MJ/kg}）を燃やすなら、何トン分にあたりますか。1回の分裂のエネルギーは \qty{200}{MeV} とします。
\phbamen
燃料の中のウラン235 の原子が、すべて分裂し尽くす場面です。
\phmushi
分裂しないで中性子を吸うだけのウラン235 は無視し、全部が分裂するとします。1回のエネルギーは、熱になる分の目安の \qty{200}{MeV} とします。
\phzu
モル質量 \qty{235}{g/mol}、$N_{\mathrm A}=\qty{6.02e23}{mol^{-1}}$、$\qty{1}{MeV}=\qty{1.60e-13}{J}$。
\phriyu
出るエネルギーは「原子の数×1回のエネルギー」です。原子の数は、物質量（mol）×アボガドロ定数で求めます。
\phtateru
\[ N=\frac{1}{235}\times N_{\mathrm A},\qquad E=N\times\qty{200}{MeV} \]
\phtoku
\[ N=\frac{1}{235}\times6.02\times10^{23}=2.56\times10^{21} \]
\[ E=2.56\times10^{21}\times200\times1.60\times10^{-13}=\qty{8.2e10}{J} \]
\qty{1}{MW} で1日（\qty{86400}{s}）は $\qty{8.64e10}{J}$ なので、$8.2\div8.64=0.95$ 日分です。炭素なら、$8.2\times10^{10}\div(3.28\times10^7)=\qty{2.5e3}{kg}$、約2.5トンです。
\phtashikame
単位：個×MeV×J/MeV＝J。\\
大きさ：1円玉1枚と同じ質量の燃料が、炭素2.5トン分。けたが約250万倍違うのは、①の「1回のエネルギーが約5千万倍」と、「原子1個の質量がウランのほうが約20倍重い」ことの両方からです。\\
極端な場合：半分の \qty{0.5}{g} なら、エネルギーも半分です。
\kotae 約 $\qty{8.2e10}{J}$（\qty{1}{MW} で約0.95日分）。炭素約2.5トン分
\end{phreidai}
\begin{jibun}
\qty{1}{W} の熱を出し続けるには、1秒に何回の分裂が必要ですか（1回 \qty{200}{MeV}）。
\bkjans{$1\div(200\times1.60\times10^{-13})\approx\num{3.1e10}$ 回}
\end{jibun}

\label{ph:原-15:ex2}
\begin{phreidai}[熱出力 3 GW の原子炉が1日に使うウラン235]
熱出力 \qty{3.0}{GW}（電気出力約 \qty{1}{GW}）の原子炉で、1秒あたり・1日あたりの分裂の回数と、1日に分裂するウラン235 の質量を求めましょう。1回は \qty{200}{MeV} とします。
\phbamen
一定の熱出力で運転している原子炉で、燃料が減っていく速さを見積もる場面です。
\phmushi
ウラン238 から生まれるプルトニウムの分裂は無視し、熱は全部ウラン235 の分裂から出るとします。
\phzu
$P=\qty{3.0e9}{W}$、1回 $\qty{200}{MeV}=\qty{3.20e-11}{J}$。
\phriyu
「熱出力÷1回のエネルギー」が1秒の分裂の回数です。回数を mol にすれば、質量がわかります（例題1 の逆向き）。
\phtateru
\[ \text{1秒の回数}=\frac{P}{\qty{3.20e-11}{J}},\qquad m=\frac{\text{1日の回数}}{N_{\mathrm A}}\times\qty{235}{g/mol} \]
\phtoku
\[ \frac{3.0\times10^9}{3.20\times10^{-11}}=9.4\times10^{19}\ \text{回/s},\qquad 9.4\times10^{19}\times\num{86400}=8.1\times10^{24}\ \text{回/日} \]
\[ m=\frac{8.1\times10^{24}}{6.02\times10^{23}}\times235\approx\qty{3.2e3}{g} \]
\phtashikame
単位：W÷J＝1/s。\\
大きさ：1日に約 \qty{3.2}{kg}、1年で約1.2トン。実際には、分裂せずに中性子を吸うだけのウラン235 もあるので（$\sigma_{\mathrm c5}/\sigma_{\mathrm f5}=0.17$）、減るウラン235 は約1.17倍の約 \qty{3.7}{kg} です。\\
極端な場合：熱出力が3分の1なら、すべて3分の1。
\kotae 1秒に約 $\num{9.4e19}$ 回、1日に約 $\num{8.1e24}$ 回、約 \qty{3.2}{kg}
\end{phreidai}
\begin{jibun}
熱出力 \qty{1.0}{GW} の原子炉では、1日に分裂するウラン235 は何 g ですか。
\bkjans{1秒に $1.0\times10^9\div(3.20\times10^{-11})=\num{3.1e19}$ 回、1日に $\num{2.7e24}$ 回。$2.7\times10^{24}\div(6.02\times10^{23})=\qty{4.5}{mol}$、$4.5\times235\approx\qty{1.1e3}{g}$（約 \qty{1.1}{kg}）}
\end{jibun}

\label{ph:原-15:ex3}
\begin{phreidai}[増倍率と世代の数]
はじめに中性子が1000個あります。（1）増倍率が $k=2$ と $k=1$ のとき、10世代あとの中性子はそれぞれ何個ですか。（2）1世代を $\qty{1e-4}{s}$ とし、遅発中性子がないとすると、$k=1.001$ で1秒後には何倍になりますか。
\phbamen
連鎖反応で、世代ごとに中性子の数が変わっていく場面です。
\phmushi
世代は全部同じ長さとし、$k$ は一定とします。遅発中性子は考えません（⑥でもどります）。
\phzu
$n_0=1000$、世代の数 $g$。
\phriyu
1世代ごとに $k$ 倍になるので、$n_g=n_0k^g$（②）を使います。
\phtateru
\[ n_{10}=1000\times k^{10},\qquad g=\frac{\qty{1}{s}}{\qty{1e-4}{s}}=\num{10000} \]
\phtoku
（1）$k=2$：$1000\times2^{10}=1000\times1024=\num{1024000}$ 個。$k=1$：$1000\times1^{10}=1000$ 個。\\
（2）$1.001^{10000}\approx\num{2.2e4}$ 倍（電卓で $e^{10000\times\ln1.001}=e^{9.995}$）。
\phtashikame
単位：個×（単位のない数）＝個。\\
大きさ：$k=2$ の10世代は、1世代 $10^{-4}$ 秒なら1000分の1秒で約1000倍です。（2）は $e^{(k-1)t/\ell}=e^{10}=\num{2.2e4}$ と同じで、④の式とも合います。\\
極端な場合：$k=1$ はどれだけ世代を重ねても変わりません。
\kotae （1）約 $\num{1.0e6}$ 個と1000個（2）約 $\num{2.2e4}$ 倍
\end{phreidai}
\begin{jibun}
はじめ1000個の中性子が、$k=0.9$ と $k=1.1$ で10世代たつと、それぞれ約何個ですか。
\bkjans{$1000\times0.9^{10}\approx349$ 個、$1000\times1.1^{10}\approx\num{2594}$ 個（表\ref{ph:原-15:tab-sheet}と同じ）}
\end{jibun}

\label{ph:原-15:ex4}
\begin{phreidai}[濃縮度で $\eta$ が変わる]
ウラン235 の割合（濃縮度）が 0.72\,\%（天然ウラン）・3\,\%・5\,\%の燃料について、燃料に熱中性子が1個吸われたときに生まれる中性子の数 $\eta$ を求めましょう。$\nu=2.43$、$\sigma_{\mathrm f5}=\qty{585}{b}$、$\sigma_{\mathrm c5}=\qty{99}{b}$、$\sigma_{\mathrm c8}=\qty{2.68}{b}$ とします。
\phbamen
燃料の中で、熱中性子がウラン235 とウラン238 のどちらに、どう吸われるかを数える場面です。
\phmushi
燃料の酸素や、ほかの物質に吸われる分は $f$ に入れるので、ここでは考えません。断面積は熱中性子1つの値で代表させます。
\phzu
$N_5$：$N_8$＝濃縮度：（1－濃縮度）。
\phriyu
吸われる確率は「数×断面積」に比例するので、⑤の $\eta$ の式を使います。
\phtateru
\[ \eta=2.43\times\frac{585N_5}{684N_5+2.68N_8} \]
\phtoku
\[ 0.72\,\%：\ \eta=2.43\times\frac{585\times0.0072}{684\times0.0072+2.68\times0.9928} \]
\[ =2.43\times\frac{4.21}{4.93+2.66}=1.35 \]
\[ 3\,\%：\ \eta=2.43\times\frac{17.55}{20.52+2.60}=1.85,\qquad 5\,\%：\ \eta=2.43\times\frac{29.25}{34.2+2.55}=1.93 \]
\phtashikame
単位：b÷b で単位なし。\\
大きさ：どれも $\nu=2.43$ より小さい。天然ウランでは、吸われる中性子の3分の1ほどがウラン238 に取られます。\\
極端な場合：ウラン238 がない（100\,\%）と、$\eta=2.43\times585/684=2.08$ が上限です。
\kotae 1.35・1.85・1.93
\end{phreidai}
\begin{jibun}
ウラン235 が 100\,\%の燃料の $\eta$ を求め、3\,\%のときの何倍かを答えましょう。
\bkjans{$2.43\times585\div684\approx2.08$、$2.08\div1.85\approx1.13$ 倍（濃縮度は約33倍なのに、$\eta$ は1割ほどしか増えない）}
\end{jibun}

\label{ph:原-15:ex5}
\begin{phreidai}[濃縮には何段いる？]
熱-14 の例題12 では、六フッ化ウランの気体を小さな穴に1回通すと、ウラン235 とウラン238 の数の比が1.0043倍になりました。天然ウラン（ウラン235 が 0.7204\,\%、ウラン238 が 99.2742\,\%）を、ウラン235 が 3\,\%の燃料にするには、理想的には何段通せばよいですか。
\phbamen
わずかな違いを何回も積み重ねて、ウランを濃縮する場面です。
\phmushi
1段ごとに比がちょうど1.0043倍になり、量が減ることや、もどってくる流れは考えません（理想の段）。
\phzu
はじめの比 $r_0=0.7204/99.2742$、目標の比 $r=3/97$、段の数 $m$。
\phriyu
1段ごとに同じ倍率でかけ算されるので、$r=r_0\times1.0043^m$ の形です。$m$ は対数で求めます（3-12）。
\phtateru
\[ 1.0043^m=\frac{r}{r_0},\qquad m=\frac{\ln(r/r_0)}{\ln1.0043} \]
\phtoku
\[ r_0=0.007257,\qquad r=0.03093,\qquad \frac{r}{r_0}=4.26 \]
\[ m=\frac{\ln4.26}{\ln1.0043}=\frac{1.450}{0.00429}\approx338 \]
\phtashikame
単位：比÷比で、段の数は単位なし。\\
大きさ：約340段。実際の装置では1段の分け方が理想より悪いので、もっと多くの段が必要になります。\\
極端な場合：比を2倍にするだけでも $\ln2/\ln1.0043\approx162$ 段かかります。
\kotae 約340段（理想）
\end{phreidai}
\begin{jibun}
同じ方法で、天然ウランをウラン235 が 5\,\%の燃料にするには、理想的に何段いりますか。
\bkjans{$r=5/95=0.05263$、$r/r_0=7.25$、$\ln7.25\div\ln1.0043\approx460$ 段}
\end{jibun}

\label{ph:原-15:ex6}
\begin{phreidai}[4因子公式で $k_\infty$ を求める]
ウラン235 が 3\,\%の燃料と水でできた、大きな原子炉の $k_\infty$ を求めましょう。$\eta=1.845$（例題4）、$\varepsilon=1.03$、$p=0.75$、$f=0.90$ とします。
\phbamen
中性子の一生を4つの段に分け、1世代でどれだけ増えるかを数える場面です（図\ref{ph:原-15:fig-life}）。
\phmushi
炉は十分に大きく、外へもれる中性子はないとします。
\phzu
$\eta=1.845$、$\varepsilon=1.03$、$p=0.75$、$f=0.90$。
\phriyu
4つの段は順に起こるので、割合のかけ算（⑤の4因子公式）を使います。
\phtateru
\[ k_\infty=\eta\,\varepsilon\,p\,f \]
\phtoku
\[ k_\infty=1.845\times1.03\times0.75\times0.90=1.28 \]
\phtashikame
単位：どれも単位なし。\\
大きさ：1より大きいので、このままでは中性子が増えていきます。実際の炉では、外へもれる分と、制御棒などに吸わせる分で、生まれる中性子の約22\,\%（$1-1/1.28$）を差し引いて、$k=1$ にしています。\\
極端な場合：天然ウラン（$\eta=1.35$）で、ほかの3つが同じなら $1.35\times1.03\times0.75\times0.90=0.94$ で、1にとどきません。
\kotae $k_\infty\approx1.28$
\end{phreidai}
\begin{jibun}
例題6 で、水などに吸われる熱中性子が増えて $f=0.85$ になると、$k_\infty$ はいくらになりますか。
\bkjans{$1.845\times1.03\times0.75\times0.85\approx1.21$}
\end{jibun}

\label{ph:原-15:ex7}
\begin{phreidai}[はだかの炉が臨界になる大きさ]
$k_\infty=1.10$、移動面積 $M^2=\qty{60}{cm^2}$ の一様な物質で、まわりに反射体のない炉を作ります。臨界になる平板の厚さ、球の半径、立方体の1辺を求めましょう。球と立方体では、どちらの体積が小さいですか。
\phbamen
炉心の形と大きさを、もれる中性子を考えて決める場面です（⑦）。
\phmushi
1群の拡散の式で考え、炉心のはしで中性子束が0になるとします（はしの外への少しのはみ出しは無視）。
\phzu
$k_\infty=1.10$、$M^2=60$、平板の厚さ $a$、球の半径 $R$、立方体の1辺 $a$。
\phriyu
臨界の条件は $B^2=(k_\infty-1)/M^2$ です。$B$ と大きさの関係は、形ごとに⑦で求めました。
\phtateru
\[ B=\sqrt{\frac{k_\infty-1}{M^2}},\qquad \text{平板}\ a=\frac{\pi}{B},\quad \text{球}\ R=\frac{\pi}{B},\quad \text{立方体}\ a=\frac{\sqrt3\,\pi}{B} \]
\phtoku
\[ B=\sqrt{\frac{0.10}{60}}=\qty{0.0408}{cm^{-1}},\qquad \frac{\pi}{B}=\qty{77}{cm},\qquad \frac{\sqrt3\,\pi}{B}=\qty{133}{cm} \]
球の体積 $\tfrac43\pi\times77^3\approx\qty{1.9e6}{cm^3}$（\qty{1.9}{m^3}）、立方体 $133^3\approx\qty{2.4e6}{cm^3}$（\qty{2.4}{m^3}）。
\phtashikame
単位：$\sqrt{1/\mathrm{cm^2}}$＝1/cm、$\pi/B$ は cm。\\
大きさ：⑦のプログラムでも、厚さ \qty{77}{cm} の平板で $k=1.000$ でした。\\
極端な場合：$k_\infty$ が1に近づくと $B\to0$ で、いくら大きくしても臨界になりにくくなります。同じ体積なら表面がいちばん小さい球が、もれがいちばん少なく、体積も小さくてすみます。
\kotae 平板の厚さ約 \qty{77}{cm}、球の半径約 \qty{77}{cm}、立方体の1辺約 \qty{133}{cm}。球のほうが体積が小さい
\end{phreidai}
\begin{jibun}
例題7 で $k_\infty=1.05$ にすると、臨界になる球の半径はいくらですか。
\bkjans{$B=\sqrt{0.05/60}=\qty{0.0289}{cm^{-1}}$、$R=\pi/B\approx\qty{109}{cm}$}
\end{jibun}

\label{ph:原-15:ex8}
\begin{phreidai}[反応度と炉の周期]
増倍率が $k=1.001$ と $k=1.003$ になったとき、反応度 $\rho$ と、それが何ドルかを求めましょう。また、即発中性子だけのとき（$\ell=\qty{1e-4}{s}$）の炉の周期 $\ell/(k-1)$ と、遅発中性子を1群で入れたときの周期 $(\beta-\rho)/(\lambda\rho)$ を比べましょう（$\beta=0.0065$、$\lambda=\qty{0.08}{s^{-1}}$）。
\phbamen
制御棒を少し引きぬいて、出力を上げはじめる場面です。
\phmushi
温度による反応度の変化（⑩）は無視し、反応度は一定とします。周期は、即発跳躍が終わったあとのゆっくりした増え方で考えます。
\phzu
$k$、$\rho=(k-1)/k$、$\beta=0.0065$、$\lambda=0.08$、$\ell=\qty{1e-4}{s}$。
\phriyu
反応度の定義（⑥）と、即発だけの式（④）、遅発中性子を入れた式（⑧の $1/s_1$）を使います。
\phtateru
\[ \rho=\frac{k-1}{k},\qquad \frac{\rho}{\beta}\ \text{ドル},\qquad T_{\text{即発}}=\frac{\ell}{k-1},\qquad T=\frac{\beta-\rho}{\lambda\rho} \]
\phtoku
$k=1.001$：$\rho=0.001/1.001=0.000999$（0.15ドル）、$T_{\text{即発}}=0.0001\div0.001=\qty{0.1}{s}$、$T=(0.0065-0.001)\div(0.08\times0.001)\approx\qty{69}{s}$。\\
$k=1.003$：$\rho=0.00299$（0.46ドル）、$T_{\text{即発}}=\qty{0.033}{s}$、$T=(0.0065-0.003)\div(0.08\times0.003)\approx\qty{15}{s}$。
\phtashikame
単位：$\ell$ は s、$\lambda$ は 1/s なので、どちらの周期も s。\\
大きさ：遅発中性子があると、周期は数百倍長くなります。\\
極端な場合：$\rho$ が $\beta$（1ドル）に近づくと $T\to0$ で、遅発中性子の効き目がなくなります。
\kotae $k=1.001$：$\rho\approx0.0010$（0.15ドル）、0.1 s と約69 s。$k=1.003$：$\rho\approx0.0030$（0.46ドル）、0.033 s と約15 s
\end{phreidai}
\begin{jibun}
$k=1.002$ のとき、反応度・ドル・即発だけの周期・遅発中性子1群の周期を求めましょう。
\bkjans{$\rho\approx0.0020$、0.31ドル、\qty{0.05}{s}、$T=0.0045\div0.00016\approx\qty{28}{s}$}
\end{jibun}

\label{ph:原-15:ex9}
\begin{phreidai}[即発跳躍と、60秒後の出力]
1点炉の動特性の式（⑧）で、反応度 $\rho=0.001$ を階段のように入れます。即発跳躍の大きさと、60秒後の出力（はじめの何倍か）を、2つの指数の近似の式で求め、⑧のルンゲ・クッタ法の結果（2.823倍）と比べましょう。$\beta=0.0065$、$\lambda=\qty{0.08}{s^{-1}}$ とします。
\phbamen
制御棒を少し引きぬいた直後から、出力がどう増えるかを、式で見積もる場面です。
\phmushi
$s_2$ の項（約0.02秒で消える）は、跳びの大きさとしてだけ考えます。温度によるフィードバックは考えません。
\phzu
はじめの出力 $n_0=1$、即発跳躍のあとの出力 $n_{\text{跳}}$、ゆっくりした増え方 $e^{s_1t}$。
\phriyu
⑧で、跳びは $\beta/(\beta-\rho)$、そのあとは $s_1\approx\lambda\rho/(\beta-\rho)$ で増えると導きました。
\phtateru
\[ n(t)\approx\frac{\beta}{\beta-\rho}\,e^{s_1t},\qquad s_1=\frac{\lambda\rho}{\beta-\rho} \]
\phtoku
\[ \frac{\beta}{\beta-\rho}=\frac{0.0065}{0.0055}=1.18,\qquad s_1=\frac{0.08\times0.001}{0.0055}=\qty{0.0145}{s^{-1}} \]
\[ n(60)\approx1.18\times e^{0.0145\times60}=1.18\times2.39=2.83 \]
\phtashikame
単位：$s_1$ は 1/s で、$s_1t$ は単位なし。\\
大きさ：ルンゲ・クッタ法の2.823倍と、3けためまでほぼ一致します。\\
極端な場合：$\rho=0$ なら跳びは1倍、$s_1=0$ で、出力は変わりません。
\kotae 跳び約1.18倍、60秒後に約2.8倍（数値の計算と合う）
\end{phreidai}
\begin{jibun}
例題9 で $\rho=0.002$ のとき、即発跳躍の大きさと60秒後の出力を、近似の式で求めましょう（ルンゲ・クッタ法では12.1倍）。
\bkjans{$0.0065\div0.0045\approx1.44$ 倍、$s_1=0.08\times0.002\div0.0045=\qty{0.0356}{s^{-1}}$、$1.44\times e^{2.13}\approx12$ 倍}
\end{jibun}

\label{ph:原-15:ex10}
\begin{phreidai}[温度が上がると反応度が下がる]
燃料と水を合わせた温度係数が $-\qty{3}{pcm/K}$（$-3\times10^{-5}$ /K）の原子炉があります。（1）運転を始めるとき、温度が \qty{300}{K} 上がると、$k$ はいくら変わりますか。（2）反応度0.001を入れたとき、温度が何 K 上がると打ち消されますか。
\phbamen
出力が上がって温度が上がり、負のフィードバックがはたらく場面です（⑨⑩）。
\phmushi
温度係数は温度によらず一定とします（実際は温度で少し変わります）。
\phzu
温度係数 $\alpha_T=-3\times10^{-5}$ /K（例として置いた値）、温度の上がり $\Delta\theta$。
\phriyu
反応度の変化は「温度係数×温度の上がり」です（反応度係数の定義）。
\phtateru
\[ \Delta k\approx\alpha_T\,\Delta\theta,\qquad \Delta\theta=\frac{0.001}{|\alpha_T|} \]
\phtoku
（1）$\Delta k=-3\times10^{-5}\times300=-0.009$。（2）$\Delta\theta=0.001\div(3\times10^{-5})\approx\qty{33}{K}$。
\phtashikame
単位：1/K×K＝単位なし。\\
大きさ：（1）の $-0.009$ は、反応度0.001の9倍で、遅発中性子の割合 $\beta=0.0065$ よりも大きい量です。冷えた状態から熱い運転の状態にするまでに、制御棒をその分だけ引きぬいておく必要があります。\\
極端な場合：温度係数が正だと、温度が上がるほど $k$ が上がり、出力がさらに上がってしまいます。軽水炉は、そうならないように作られています。
\kotae （1）$\Delta k\approx-0.009$（2）約 \qty{33}{K}
\end{phreidai}
\begin{jibun}
温度係数が $-\qty{2}{pcm/K}$ の原子炉に、反応度0.0015を入れました。温度が何 K 上がると打ち消されますか。
\bkjans{$0.0015\div(2\times10^{-5})=\qty{75}{K}$}
\end{jibun}

\label{ph:原-15:ex11}
\begin{phreidai}[止めたあとのキセノン]
⑩のキセノン135 の式で、運転中のつり合い（$\dd I/\dd t=\dd X/\dd t=0$）から、ヨウ素135 とキセノン135 の数の比 $I_0/X_0$ を求めましょう。次に、止めてから8時間後のキセノン135 の量が、運転中の何倍かを求めましょう。$\gamma_I=0.064$、$\gamma_X=0.0024$、$\sigma_X=\qty{2.6e6}{b}$、中性子束 $\phi=\qty{3e13}{cm^{-2}.s^{-1}}$、半減期はヨウ素135 が \qty{6.58}{h}、キセノン135 が \qty{9.14}{h} とします。
\phbamen
運転を止めたあとに、中性子を強く吸うキセノン135 がたまる場面です。
\phmushi
運転中の中性子束は一定、止めた瞬間に0になるとします。ヨウ素135 が中性子を吸う分は無視します。
\phzu
\[ \lambda_I=\frac{\ln2}{6.58}=\qty{0.1053}{h^{-1}},\qquad \lambda_X=\frac{\ln2}{9.14}=\qty{0.0758}{h^{-1}} \]
\[ \sigma_X\phi=2.6\times10^{-18}\times3\times10^{13}\times3600=\qty{0.281}{h^{-1}} \]
（断面積は $\qty{1}{b}=\qty{e-24}{cm^2}$ で cm$^2$ に直し、3600 をかけて 1 時間あたりにした）。
\phriyu
つり合いの式で $I_0$・$X_0$ を求めます。止めたあとは $\phi=0$ で、原-12 ⑧の親と子の式（6-14）になります。
\phtateru
\[ I_0=\frac{\gamma_I\Sigma_{\mathrm f}\phi}{\lambda_I},\qquad X_0=\frac{(\gamma_I+\gamma_X)\Sigma_{\mathrm f}\phi}{\lambda_X+\sigma_X\phi} \]
\[ \frac{X}{X_0}=e^{-\lambda_Xt}+\frac{I_0}{X_0}\,\frac{\lambda_I}{\lambda_X-\lambda_I}\bigl(e^{-\lambda_It}-e^{-\lambda_Xt}\bigr) \]
\phtoku
\[ \frac{I_0}{X_0}=\frac{0.064}{0.0664}\times\frac{0.0758+0.281}{0.1053}=3.26 \]
$t=8$ h：$e^{-0.607}=0.545$、$e^{-0.843}=0.430$ なので、
\[ \frac{X}{X_0}=0.545+3.26\times\frac{0.1053}{-0.0295}\times(0.430-0.545)=0.545+1.34=1.88 \]
\phtashikame
単位：/h×h で指数の肩は単位なし。\\
大きさ：止めて8時間後に、運転中の約1.9倍。図\ref{ph:原-15:fig-xenon}の山の高さ1.88と同じです。4時間後は1.70倍、24時間後は1.12倍、48時間後は0.26倍です。\\
極端な場合：中性子束が小さい炉では、運転中も崩壊で消える分が多く、$I_0/X_0$ が小さいので、山は低くなります。
\kotae $I_0/X_0\approx3.3$、8時間後に約1.9倍
\end{phreidai}
\begin{jibun}
例題11 で、中性子束が $\qty{1e13}{cm^{-2}.s^{-1}}$ の炉では、$I_0/X_0$ はいくらですか。止めたあとのキセノンの山は、例題11 より高いですか、低いですか。
\bkjans{$\sigma_X\phi=\qty{0.0936}{h^{-1}}$、$0.0758+0.0936=0.169$、$I_0/X_0=0.964\times0.169\div0.1053$ $\approx1.55$。山は低い（計算すると約1.2倍）}
\end{jibun}

\label{ph:原-15:ex12}
\begin{phreidai}[止めたあとの崩壊熱]
熱出力 \qty{3.0}{GW} で1年（$T_0=\qty{3.16e7}{s}$）運転した原子炉を止めました。止めてから1秒後・1時間後・1日後・1年後の崩壊熱を、⑪の目安の式で求めましょう。
\phbamen
連鎖反応を止めたあとも出続ける熱の大きさを見積もり、冷やすための設備の大きさを考える場面です。
\phmushi
何百種類もの分裂片の崩壊を、1つの目安の式にまとめます（精密な値とは2倍ほど違うこともある）。運転中の出力は一定とします。
\phzu
$P_0=\qty{3.0}{GW}$、$T_0=\qty{3.16e7}{s}$、止めてからの時間 $t$（s）。
\phriyu
止めたあとの熱は分裂片の崩壊なので、⑪の崩壊熱の式を使います。
\phtateru
\[ \frac{P}{P_0}=0.066\left[t^{-0.2}-(t+T_0)^{-0.2}\right] \]
\phtoku
$(3.16\times10^7)^{-0.2}=0.0316$ なので、\\
1秒後：$0.066\times(1-0.0316)=0.0639$（6.4\,\%、\qty{190}{MW}）\\
1時間後（\qty{3600}{s}）：$0.066\times(0.1944-0.0316)=0.0107$（1.1\,\%、\qty{32}{MW}）\\
1日後（\qty{86400}{s}）：$0.066\times(0.1030-0.0316)=0.0047$（0.47\,\%、\qty{14}{MW}）\\
1年後：$0.066\times(0.0316-0.0275)=0.00027$（0.027\,\%、\qty{0.8}{MW}）
\phtashikame
単位：両辺とも単位なし（$t$ は s で入れる約束）。\\
大きさ：1日後の \qty{14}{MW} は、100\,℃の水（蒸発熱 \qty{2257}{kJ/kg}）を1秒に約 \qty{6}{kg}、1時間に約22トン蒸発させる熱です。水を入れ続けないと、炉心の水はどんどん減ります。\\
極端な場合：$t\to0$ でも式は有限（約6.6\,\%）で、運転中の熱の約7\,\%が崩壊熱だったこと（①）と合います。
\kotae 約 \qty{190}{MW}・\qty{32}{MW}・\qty{14}{MW}・\qty{0.8}{MW}
\end{phreidai}
\begin{jibun}
例題12 の原子炉で、止めてから10日後（\qty{864000}{s}）の崩壊熱は約何 MW ですか（$\num{864000}^{-0.2}=0.0650$、$(t+T_0)^{-0.2}=0.0315$）。
\bkjans{$0.066\times(0.0650-0.0315)\approx0.0022$、$3.0\times10^9\times0.0022\approx\qty{6.6}{MW}$}
\end{jibun}

\label{ph:原-15:ex13}
\begin{phreidai}[D-T 核融合の燃料 1 g]
重水素（\qty{2.0141}{u}）と三重水素（\qty{3.0160}{u}）を同じ数ずつ混ぜた燃料 \qty{1}{g} が、全部 D＋T → $^4$He＋n の反応をしたら、何 J のエネルギーが出ますか。ウラン235 の \qty{1}{g} の分裂（例題1）の何倍ですか。1回のエネルギーは \qty{17.6}{MeV} とします。
\phbamen
核融合の燃料が持つエネルギーを、核分裂と比べる場面です（⑫）。
\phmushi
燃料が全部反応するとします（実際には、ごく一部しか反応しない）。中性子が持ち去るエネルギーも全部数えます。
\phzu
DとTの組の質量 $2.0141+3.0160=\qty{5.030}{u}$、つまり \qty{5.030}{g/mol}。
\phriyu
例題1 と同じく「反応の回数×1回のエネルギー」です。1回の反応には、DとTが1個ずつ使われます。
\phtateru
\[ \text{組の数}=\frac{1}{5.030}\times N_{\mathrm A},\qquad E=\text{組の数}\times\qty{17.6}{MeV} \]
\phtoku
\[ \frac{6.02\times10^{23}}{5.030}=1.20\times10^{23} \]
\[ E=1.20\times10^{23}\times17.6\times1.60\times10^{-13}=\qty{3.4e11}{J} \]
例題1 の $\qty{8.2e10}{J}$ の約4.1倍です。
\phtashikame
単位：個×MeV×J/MeV＝J。\\
大きさ：1回のエネルギーは核分裂の約11分の1ですが、燃料の原子がずっと軽いので、質量あたりでは約4倍になります（1 u あたり、D-T は約 \qty{3.5}{MeV}、核分裂は約 \qty{0.85}{MeV}）。\\
極端な場合：燃料の半分が反応しなければ、出るエネルギーも半分です。
\kotae 約 $\qty{3.4e11}{J}$、ウラン235 の分裂の約4倍
\end{phreidai}
\begin{jibun}
ITER がめざす核融合の出力 \qty{500}{MW} を出すには、1秒に何回の D-T 反応が必要ですか。1秒に使う燃料は何 mg ですか。
\bkjans{1回は $17.6\times1.60\times10^{-13}=\qty{2.82e-12}{J}$、$5.00\times10^8\div(2.82\times10^{-12})$ $\approx\num{1.8e20}$ 回。$1.8\times10^{20}\div(6.02\times10^{23})=\qty{3.0e-4}{mol}$、$\times\,\qty{5.03}{g/mol}\approx\qty{1.5}{mg}$}
\end{jibun}

%% ======================== まちがえやすいところ ========================
\Needspace{11\baselineskip}\bksection{tsumazuki}{まちがえやすいところをチェック}
\begin{tsumazuki}{「臨界」は爆発の一歩手前？}
\bkng{原子炉が臨界になったというのは、危ない状態になったということ。}
\bkok{臨界は $k=1$、中性子の数がつり合って、一定の出力で落ち着いている状態。運転中の原子炉はいつも臨界。}
物理の「臨界」は、増えも減りもしない状態の名前です（②④）。
\end{tsumazuki}

\begin{tsumazuki}{制御棒を入れれば、熱はすぐに止まる？}
\bkng{制御棒を全部入れて連鎖反応を止めれば、原子炉は冷める。}
\bkok{連鎖反応は止まるが、分裂片の崩壊熱が出続ける。1日後でも運転中の約0.5\,\%。冷やし続けないといけない。}
崩壊熱は、時間の $-0.2$ 乗でしか減りません（⑪）。
\end{tsumazuki}

\begin{tsumazuki}{$k=1.001$ なら、ほとんど変わらない？}
\bkng{1.001 は1とほとんど同じだから、出力もほとんど変わらない。}
\bkok{1世代ごとに0.1\,\%ずつでも、世代はとても短いので、積み重なると大きく増える。遅発中性子があっても、約1分で $e$ 倍。}
増倍率のわずかなずれが、指数の肩にのって効いてきます（④⑥）。
\end{tsumazuki}

\begin{tsumazuki}{ウランは燃えている？}
\bkng{原子炉の中では、ウランが酸素と結びついて燃えている。}
\bkok{燃焼（化学の反応）ではなく、原子核が分かれる核分裂。酸素はいらず、1回のエネルギーは化学の反応の約5千万倍。}
「核燃料が燃える」という言い方は、たとえとして使われることばです（①）。
\end{tsumazuki}

\begin{tsumazuki}{濃縮度を2倍にすれば、$k$ も2倍？}
\bkng{ウラン235 を2倍にすれば、分裂も中性子も2倍になり、$k$ も2倍になる。}
\bkok{$\eta$ はウラン238 との取り合いで決まり、3\,\%から6\,\%で約5\,\%しか増えない（上限2.08）。}
例題4 と⑤の2倍の図を見直しましょう。
\end{tsumazuki}

\begin{column}{スカッシュのコートの下の原子炉}
1942年12月2日、アメリカのシカゴ大学のスタジアムの下にあったスカッシュのコートで、フェルミの組は、黒鉛のブロックとウランを何層にも積み上げた「パイル（積み重ね）」とよばれる装置を組み上げました。黒鉛は、中性子をほとんど吸わない減速材です。中性子を吸う棒を少しずつ引きぬいていくと、中性子の数を数える計器の針が、止まらずに増え続けるようになりました。人類が、核分裂の連鎖反応を、はじめて制御しながら続けた瞬間です。核分裂の発見から、まだ4年しかたっていませんでした。フェルミは、その4年前の1938年に、遅い中性子による核反応の研究でノーベル物理学賞を受けていました。
\end{column}

%% ======================== 工学ではこう使う ========================
\Needspace{18\baselineskip}\bksection{phgear}{工学ではこう使う}
\begin{phkougaku}[原子炉を設計する・運転する]
\textbf{炉心を設計する}：燃料の濃縮度と並べ方、燃料と水の割合を決め、$k$ と出力の形（中性子束の形：⑦）を計算します。まん中だけが熱くなりすぎないよう、濃縮度の違う燃料を組み合わせて、出力を平らに近づけます。温度係数がいつも負になること（⑩）も、設計で確かめます。

\textbf{運転する}：炉心のまわりや中に中性子の検出器（原-14）を置いて出力をはかり、制御棒の位置などで $k$ を調節します。燃料が減り、分裂片がたまると $k$ が下がっていくので、運転の期間全体を通して $k$ を1に保てるように、はじめに余分な反応度を用意しておきます（例題6 の22\,\%）。

\textbf{止めたあとも冷やす}：崩壊熱（⑪）を取り除くため、原子炉には、ふだんの冷却とは別の冷やす設備が、いくつも用意されます。どんな場合にも電源と水を確保することが、設計と規制の中心になっています（原-16）。

\begin{phketa}[原子炉の数の目安]
\phketaline{1回の分裂のエネルギー}{約 \qty{200}{MeV}（化学の反応の約5千万倍）}
\phketaline{熱出力 3 GW の炉}{1日にウラン235 約 \qty{3}{kg} を分裂させる}
\phketaline{遅発中性子の割合}{約0.65\,\%（平均約13秒遅れる）}
\phketaline{中性子の1世代}{$10^{-4}$ 秒ほど（遅発中性子をふくむ平均は約0.08秒）}
\phketaline{崩壊熱}{止めた直後 約6〜7\,\%、1日後 約0.5\,\%}
\phboxhead{答えが桁で合っているか}
出力から分裂の回数を出すときは、「\qty{1}{W} は1秒に約 $3\times10^{10}$ 回」を覚えておくと、けたの見当がつきます。GW なら $3\times10^{19}$ 回/s です。$10^{40}$ 回のような答えが出たら、MeV と J の換算（$\qty{1}{MeV}=\qty{1.6e-13}{J}$）を確かめましょう。
\end{phketa}
\end{phkougaku}
\phstation{原子炉物理,原子炉工学,放射線工学,核融合工学}{核分裂、増倍率、臨界、反応度、遅発中性子、中性子束、拡散方程式、バックリング、反応度係数、崩壊熱}

%% ======================== 章末問題 ========================
\bkendproblems
\bkproblemgroup{基本}{3}
\bkproblem 熱出力 \qty{1.0}{GW} の原子炉を1年（365日）運転します。1回の分裂を \qty{200}{MeV} として、1年の分裂の回数と、分裂するウラン235 の質量を求めましょう。
\bkpans{1秒に $\num{3.1e19}$ 回、1年に約 $\num{9.8e26}$ 回。$9.8\times10^{26}\div(6.02\times10^{23})=\qty{1.6e3}{mol}$、$1.6\times10^3\times\qty{235}{g}$ $\approx\qty{3.8e2}{kg}$}
\bkproblem はじめ1000個の中性子が、増倍率 $k=1.05$ と $k=0.95$ で20世代たつと、それぞれ約何個ですか。
\bkpans{$1000\times1.05^{20}\approx\num{2.7e3}$ 個（2653）、$1000\times0.95^{20}\approx\num{3.6e2}$ 個（358）}
\bkproblem 増倍率が $k=1.0025$ になりました。反応度 $\rho$ と、それが何ドルかを求めましょう（$\beta=0.0065$）。
\bkpans{$\rho=0.0025\div1.0025\approx0.0025$、$0.0025\div0.0065\approx0.38$ ドル}
\bkproblem ある原子炉で、$\eta=2.0$、$\varepsilon=1.02$、$p=0.80$、$f=0.85$ です。$k_\infty$ を求めましょう。この炉で $k=1$ にするには、生まれる中性子の約何\,\%を、もれや制御棒で失わせる必要がありますか。
\bkpans{$k_\infty=2.0\times1.02$ $\times0.80\times0.85\approx1.39$。$1-1/1.39\approx0.28$、約28\,\%}
\bkproblem 熱出力 \qty{2.0}{GW} で1年運転した原子炉を止めました。⑪の目安の式で、1時間後と1日後の崩壊熱を求めましょう（$(t+T_0)^{-0.2}$ はどちらも $0.0316$、$3600^{-0.2}=0.1944$、$\num{86400}^{-0.2}=0.1030$）。
\bkpans{1時間後 $0.066\times0.1628\times2.0\,\mathrm{GW}\approx\qty{21}{MW}$、1日後 $0.066\times0.0714\times2.0\,\mathrm{GW}\approx\qty{9.4}{MW}$}
\bkproblem D＋T の核融合（1回 \qty{17.6}{MeV}、燃料の組 \qty{5.03}{u}）と、ウラン235 の分裂（1回 \qty{200}{MeV}、ウラン235 と中性子で \qty{236}{u}）を、（1）1回の反応のエネルギーと、（2）反応する原子核の質量 \qty{1}{u} あたりのエネルギーで比べましょう。
\bkpans{（1）分裂が約11倍（$200\div17.6$）。（2）D-T は $17.6\div5.03\approx\qty{3.5}{MeV}$、分裂は $200\div236\approx\qty{0.85}{MeV}$ で、D-T が約4.1倍}
\bkproblemgroup{チャレンジ}{5}
\bkproblem $k_\infty=1.08$、$M^2=\qty{50}{cm^2}$ の物質で、はだかの炉を作ります。（1）臨界になる球の半径と、立方体の1辺を求めましょう。（2）立方体の体積は、球の体積の何倍ですか。
\bkpans{（1）$B=\sqrt{0.08/50}=\qty{0.0400}{cm^{-1}}$、球 $R=\pi/B\approx\qty{79}{cm}$、立方体 $\sqrt3\pi/B\approx\qty{136}{cm}$。（2）$136^3\div(\tfrac43\pi\times79^3)\approx1.24$ 倍}
\bkproblem 1点炉（$\beta=0.0065$、$\lambda=\qty{0.08}{s^{-1}}$）に、反応度0.2ドルを階段のように入れます。（1）反応度 $\rho$ と、即発跳躍の大きさを求めましょう。（2）跳んだあとの周期と、60秒後の出力（はじめの何倍か）を、近似の式で求めましょう。
\bkpans{（1）$\rho=0.2\times0.0065=0.0013$、$0.0065\div0.0052=1.25$ 倍。（2）$T=0.0052\div0.000104=\qty{50}{s}$、$1.25\times e^{60/50}\approx4.2$ 倍（ルンゲ・クッタ法では4.14倍）}
\bkproblem 熱出力 \qty{3.0}{GW} で1年運転した原子炉を止めました（$T_0=\qty{3.16e7}{s}$）。（1）1日後の崩壊熱は、10日後の何倍ですか（$\num{86400}^{-0.2}=0.1030$、$\num{864000}^{-0.2}=0.0650$。$(t+T_0)^{-0.2}$ は1日後 0.0316、10日後 0.0315）。（2）30日後（$\qty{2.59e6}{s}$、$t^{-0.2}=0.0522$、$(t+T_0)^{-0.2}=0.0311$）の崩壊熱は何 MW ですか。止めた日から30日で、崩壊熱が「半分の半分の…」と半減期のように減らないのはなぜか、説明しましょう。
\bkpans{（1）$0.1030-0.0316=0.0714$、$0.0650-0.0315=0.0335$、$0.0714\div0.0335\approx2.1$ 倍。（2）$0.0522-0.0311=0.0211$、$0.066\times0.0211\times\qty{3.0}{GW}$ $\approx\qty{4.2}{MW}$。半減期の違う何百種類もの分裂片の崩壊の和なので、短い半減期のものが先に減ったあとは、長い半減期のものが残り、減り方がだんだんゆっくりになるから}

\begin{furikaeri}
\begin{checklist}
\item ウラン235 の核分裂で、1回に約 200 MeV のエネルギーと2〜3個の中性子が出ることと、連鎖反応・原子炉のしくみを、図で説明できる。\bkfurref{ph:原-15:sec1}
\item 増倍率 $k$・反応度・遅発中性子を使って、原子炉の出力の増え方（炉の周期）を見積もれる。\bkfurref{ph:原-15:sec6}
\item 4因子公式と1群の拡散方程式で、増倍率と、原子炉が臨界になる大きさを求められる。\bkfurref{ph:原-15:sec7}
\item 1点炉の動特性の式・反応度係数・キセノンの式を、微分方程式として読める。\bkfurref{ph:原-15:sec8}
\item 崩壊熱と核融合のエネルギーを数で見積もり、原子炉を止めたあとも冷やし続けなければならないわけを説明できる。\bkfurref{ph:原-15:sec11}
\end{checklist}
\end{furikaeri}
\bkfukidashi{sora}{原子炉が一定の出力で動けるのは、0.65\,\%の遅れてくる中性子のおかげだったんだね。止めたあとも熱が出続けるから、冷やし続けることが大事なんだ。}
